% Wysp : architecture technique, workflows et plan d'implementation
% Compilation : tectonic -X compile docs/architecture/wysp-architecture.tex
\documentclass[11pt,a4paper]{article}

\usepackage{fontspec}
\usepackage[french]{babel}
\setmainfont{texgyrepagella}[Extension=.otf, UprightFont=*-regular, BoldFont=*-bold, ItalicFont=*-italic, BoldItalicFont=*-bolditalic]
\setsansfont{texgyreheros}[Extension=.otf, UprightFont=*-regular, BoldFont=*-bold, ItalicFont=*-italic, BoldItalicFont=*-bolditalic]

\usepackage[margin=2.1cm,headheight=14pt]{geometry}
\usepackage{xcolor}
\usepackage{tikz}
\usetikzlibrary{arrows.meta,positioning,shapes.geometric,shapes.multipart,fit,backgrounds,calc,babel}
\usepackage{pgfgantt}
\usepackage{pgf-umlsd}
\usepackage{booktabs,array,multirow,xltabular}
\usepackage{enumitem}
\usepackage{listings}
\usepackage[most]{tcolorbox}
\usepackage{fancyhdr}
\usepackage{titlesec}
\usepackage{amssymb}
\usepackage[hidelinks]{hyperref}

% ---------------------------------------------------------------- couleurs
\definecolor{night}{HTML}{0F172A}
\definecolor{surface}{HTML}{1E293B}
\definecolor{wisp}{HTML}{0F9E8E}
\definecolor{wisplight}{HTML}{CCFBF1}
\definecolor{pumpkin}{HTML}{EA580C}
\definecolor{pumpkinlight}{HTML}{FFEDD5}
\definecolor{muted}{HTML}{64748B}
\definecolor{okgreen}{HTML}{16A34A}
\definecolor{danger}{HTML}{DC2626}
\definecolor{pa}{HTML}{2563EB}   % personne A
\definecolor{pb}{HTML}{9333EA}   % personne B
\definecolor{pc}{HTML}{EA580C}   % personne C
\definecolor{codebg}{HTML}{F8FAFC}

% ---------------------------------------------------------------- mise en page
\titleformat{\section}{\Large\sffamily\bfseries\color{night}}{\thesection}{0.8em}{}
\titleformat{\subsection}{\large\sffamily\bfseries\color{night}}{\thesubsection}{0.7em}{}
\titleformat{\subsubsection}{\normalsize\sffamily\bfseries\color{surface}}{\thesubsubsection}{0.6em}{}
\pagestyle{fancy}
\fancyhf{}
\fancyhead[L]{\sffamily\small\color{muted} Wysp \textbullet{} Architecture, workflows et plan}
\fancyhead[R]{\sffamily\small\color{muted} Brouillon 0.1 \textbullet{} 9 novembre 2026}
\fancyfoot[C]{\sffamily\small\thepage}
\renewcommand{\headrulewidth}{0.3pt}
\setlist{itemsep=2pt,topsep=4pt}
\setlength{\parskip}{4pt}
\setlength{\parindent}{0pt}
\setlength{\emergencystretch}{3em}
\renewcommand{\arraystretch}{1.18}

% ---------------------------------------------------------------- code
\lstdefinestyle{wysp}{
  basicstyle=\ttfamily\footnotesize, backgroundcolor=\color{codebg}, frame=single, rulecolor=\color{muted!40},
  breaklines=true, columns=fullflexible, keepspaces=true, showstringspaces=false,
  keywordstyle=\color{pa}\bfseries, commentstyle=\color{muted}\itshape, stringstyle=\color{okgreen!70!black},
  xleftmargin=2pt, xrightmargin=2pt, aboveskip=6pt, belowskip=6pt}
\lstset{style=wysp}

% ---------------------------------------------------------------- boites
\newtcolorbox{keybox}[1][]{enhanced, colback=wisplight!45, colframe=wisp, boxrule=0.6pt, arc=2mm,
  left=6pt, right=6pt, top=4pt, bottom=4pt, fonttitle=\sffamily\bfseries, title={#1}}
\newtcolorbox{warnbox}[1][]{enhanced, colback=pumpkinlight!60, colframe=pumpkin, boxrule=0.6pt, arc=2mm,
  left=6pt, right=6pt, top=4pt, bottom=4pt, fonttitle=\sffamily\bfseries, title={#1}}
\newtcolorbox{notebox}{enhanced, colback=codebg, colframe=muted!50, boxrule=0.4pt, arc=2mm,
  left=6pt, right=6pt, top=3pt, bottom=3pt}

% ---------------------------------------------------------------- macros
\newcommand{\code}[1]{\texttt{\small #1}}
\newcommand{\file}[1]{{\small\path{#1}}}
\renewcommand{\UrlFont}{\ttfamily}
\newcommand{\PA}{\textcolor{pa}{\textbf{A}}}
\newcommand{\PB}{\textcolor{pb}{\textbf{B}}}
\newcommand{\PC}{\textcolor{pc}{\textbf{C}}}
\newcommand{\pobs}{\textsf{observed}}
\newcommand{\pinf}{\textsf{inferred}}
\newcommand{\pill}{\textsf{illustrative}}
\newcolumntype{L}[1]{>{\raggedright\arraybackslash}p{#1}}
\newcolumntype{Y}{>{\raggedright\arraybackslash}X}

\begin{document}

% ---------------------------------------------------------------- titre
\begin{titlepage}
\begin{tikzpicture}[remember picture,overlay]
  \fill[night] (current page.north west) rectangle ([yshift=-9.5cm]current page.north east);
  \foreach \x/\y/\r/\o in {3.2/-2.6/0.55/0.9, 4.1/-3.4/0.25/0.6, 2.5/-3.7/0.15/0.5, 15.5/-6.8/0.35/0.5, 16.4/-7.6/0.18/0.4}
    \shade[inner color=wisp!80!white, outer color=night, opacity=\o] ([xshift=\x cm, yshift=\y cm]current page.north west) circle (\r cm);
  \node[anchor=west, text=white, font=\sffamily\bfseries\fontsize{54}{60}\selectfont] at ([xshift=2.2cm,yshift=-5.0cm]current page.north west) {Wysp};
  \node[anchor=west, text=wisplight, font=\sffamily\Large] at ([xshift=2.25cm,yshift=-6.4cm]current page.north west) {Un framework pour créer des tuteurs IA adaptés aux logiciels};
  \node[anchor=west, text=white!75, font=\sffamily\normalsize] at ([xshift=2.25cm,yshift=-7.3cm]current page.north west) {Kiroween 2026 \textbullet{} Catégorie Skeleton Crew};
\end{tikzpicture}

\vspace*{9.2cm}
{\sffamily\huge\bfseries\color{night} Architecture technique, workflows\\[2pt] et plan d'implémentation\par}
\vspace{0.6cm}
{\sffamily\large\color{muted} Brouillon de travail pour l'équipe (trois personnes)\par}
\vspace{1.0cm}

\begin{tabularx}{\textwidth}{@{}L{3.4cm}Y@{}}
\textbf{Version} & 0.1 (brouillon), 9 novembre 2026 \\
\textbf{Base} & Note de concept « Project K », renommé \textbf{Wysp} \\
\textbf{Référence étudiée} & Kiro InteractionKit, Grand Prize Kiroween 2025 \\
\textbf{Échéance} & Vendredi 4 décembre 2026, 14\,h heure du Pacifique (22\,h GMT, heure de Dakar) \\
\textbf{Dépôt} & \file{github.com/MELAS-exe/Kiro-Wysp} \\
\textbf{Mise en place Kiro} & 10 fichiers de steering, 6 specs, 10 hooks, 3 agents, configuration MCP proposée \\
\end{tabularx}

\vfill
\begin{notebox}
\small Ce document rassemble l'analyse du projet de référence, l'architecture proposée pour Wysp, les workflows principaux, la configuration Kiro mise en place dans le dépôt et le plan d'implémentation réparti entre trois personnes. Les choix techniques sont des propositions à valider pendant la première semaine. Les specs Kiro (\file{.kiro/specs/}) restent la référence détaillée, en anglais pour le jury.
\end{notebox}
\end{titlepage}

\tableofcontents
\clearpage

\section{Résumé}

\textbf{Wysp} est un socle réutilisable pour créer des tuteurs IA propres à chaque logiciel. Un tuteur Wysp observe l'état réel de l'application, explique ses concepts, construit un tableau visuel interactif, guide les actions et vérifie dans l'application ce que l'utilisateur a réussi. Le prototype doit prouver qu'un même cœur sert deux applications très différentes : \textbf{Blender} (intégration native par l'API Python) et un \textbf{tableau de bord web} (intégration par le DOM et l'accessibilité, via une extension Chrome).

\begin{keybox}[Les décisions structurantes]
\begin{enumerate}[leftmargin=1.2em]
  \item \textbf{Contrat d'abord.} Un protocole JSON-RPC 2.0 sur WebSocket local. Les adaptateurs sont clients, le cœur est serveur sur \code{127.0.0.1:8765}. Les modèles Pydantic sont la source unique, exportés en JSON Schema et en types TypeScript.
  \item \textbf{Un modèle d'observation générique} : un graphe d'entités (\code{ref}, \code{kind}, \code{props}, liens, focus). Les types d'entités sont déclarés par l'adaptateur, jamais par le cœur.
  \item \textbf{Le LLM décide quoi dire et quoi montrer. Le cœur décide ce qui est vrai et ce qui est permis.} Les nœuds « observés » du tableau sont créés par projection des observations, pas par le LLM. Les actions passent par une porte de permission, puis une vérification par prédicats, avec annulation automatique en cas d'échec.
  \item \textbf{Provenance visible partout} : \pobs, \pinf, \pill, avec le même code visuel sur le tableau et dans les sous-titres.
  \item \textbf{Vérification par preuves} : les critères de réussite sont des prédicats écrits par l'auteur du profil et évalués sur une observation fraîche. Une réussite vérifiée compte plus qu'une réponse ou une déclaration.
  \item \textbf{Testabilité} : un adaptateur factice (Kanban), un LLM scénarisé et un kit de tests de contrat exécutable contre chaque adaptateur.
\end{enumerate}
\end{keybox}

\paragraph{Ce que nous retenons du lauréat 2025.} Kiro InteractionKit a gagné avec un cœur petit et clair, deux applications distinctes, un dossier \file{.kiro} complet et un README très visuel. Leur cœur (29 fichiers, 925 lignes) ne contenait que des fonctions de détection, et chaque application était un fichier monolithique de 870 à 1\,883 lignes. Leurs 36 tâches de tests optionnelles, dont 29 tests de propriétés, n'ont pas été réalisées. Wysp reprend la forme (cœur mince, deux applications, Kiro visible, onboarding en quatre étapes) et va plus loin sur la frontière : un contrat versionné, un kit de tests de contrat et des propriétés de correction obligatoires.

\paragraph{Répartition en trois sections.}
\begin{itemize}[leftmargin=1.2em]
  \item \PA{} \textbf{Cœur et contrat} : contrat, passerelle, orchestrateur, LLM, tableau (modèle), confiance, pédagogie, kit de contrat, adaptateur factice.
  \item \PB{} \textbf{Expérience} : panneau React, serveur du panneau et narration côté cœur, rendu du tableau, voix, tableau de bord de démo, contenu pédagogique des deux profils, identité visuelle, vidéo.
  \item \PC{} \textbf{Adaptateurs} : SDK Python et TypeScript, extension Blender et scène de démo, extension Chrome.
\end{itemize}

\paragraph{Calendrier.} Il reste 25 jours. Quatre jalons : \textbf{J0} mercredi 11 novembre (contrat 1.0 figé, squelettes connectés), \textbf{J1} jeudi 19 novembre (Blender de bout en bout en mode texte), \textbf{J2} vendredi 27 novembre (navigateur de bout en bout sans toucher au cœur, voix), \textbf{J3} mardi 1\textsuperscript{er} décembre (gel des fonctionnalités). Tournage les 1\textsuperscript{er} et 2 décembre, soumission le 4 décembre avant midi GMT.

\paragraph{Mise en place Kiro faite dans le dépôt.} Dix fichiers de steering, six specs complètes (exigences EARS, design avec 49 propriétés de correction, tâches avec responsable et jalon), dix hooks au format v1 appuyés sur un script Python testé, trois agents spécialisés. La configuration MCP est fournie dans \file{docs/kiro/mcp.example.json}, à copier dans \file{.kiro/settings/mcp.json} (ce dossier est protégé en écriture pour l'agent).

\begin{warnbox}[Points à valider par l'équipe cette semaine]
Choix du fournisseur LLM et du service de synthèse vocale (budget, clés), disponibilité de Blender 5.2 sur les trois machines, répartition des rôles proposée, langue de la vidéo (anglais recommandé), et relecture des règles officielles 2026 (la page des règles n'a pas pu être lue automatiquement).
\end{warnbox}

\section{Analyse du projet de référence : Kiro InteractionKit}

Kiro InteractionKit (dépôt \file{github.com/Harpita-P/Kiro-InteractionKit}, dernier commit du 5 décembre 2025) a remporté le Grand Prize de Kiroween 2025 dans la catégorie Skeleton Crew. Nous avons cloné le dépôt et lu le code, les specs, le steering, les hooks et la configuration MCP.

\subsection{Contenu du dépôt}

\begin{xltabular}{\textwidth}{@{}L{3.6cm}L{4.4cm}Y@{}}
\toprule
\textbf{Élément} & \textbf{Contenu} & \textbf{Observation} \\
\midrule
Cœur \file{kiro_interaction_kit/} & 3 contrôleurs (main, tête, visage) sur MediaPipe, 18 fonctions de détection de gestes (\code{HA1}--\code{HA9}, \code{HE1}--\code{HE6}, \code{FA1}--\code{FA3}) & 29 fichiers, 925 lignes. Fonctions pures, une par fichier, numérotées. L'état est un dataclass de booléens (\code{is\_pinch}, \code{is\_peace}\dots). \\
Applications \file{my_apps/} & Math-O-Lantern (jeu éducatif), Holo-Board (tableau lumineux gestuel avec enregistrement vidéo) & Un seul \file{main.py} par application : 870 lignes et 7 classes, 1\,883 lignes et 5 classes. Les deux n'importent que \code{HandTracker}. \\
Specs \file{.kiro/specs/} & Une spec par application : exigences, design, tâches & Exigences en EARS (9 et 12 exigences), design avec 19 et 26 propriétés de correction, tâches cochées (13 et 12 faites). \\
Steering \file{.kiro/steering/} & \file{product}, \file{tech}, \file{structure}, \file{testing}, \file{game-development} (manuel) & Règles de nommage des gestes, séparation « vision / logique de jeu », modèle de code imposé, charte graphique des démos de test. \\
Hooks \file{.kiro/hooks/} & 3 hooks « intégrer automatiquement un nouveau geste » & À chaque nouveau fichier de geste, l'agent met à jour l'import, le champ d'état, l'appel dans le contrôleur, la démo de test et le README. Ancien format \file{.kiro.hook}. \\
MCP \file{.kiro/settings/} & Context7 & Documentation à jour de MediaPipe et Pygame pour l'agent. \\
Tests \file{tests/} & 3 démos interactives (\file{CV-Test-*.py}) & Ce sont des démos visuelles, pas des tests automatisés. \\
Onboarding & \file{setup.sh}, \file{setup.bat}, README avec GIF animés, licence MIT & « Supprimez les deux applications et créez votre spec » : le template est pensé pour être réutilisé. \\
\bottomrule
\end{xltabular}

\subsection{Leur méthode d'implémentation}

\begin{enumerate}[leftmargin=1.2em]
  \item \textbf{Un cœur qui retire une complexité répétitive.} Le kit cache MediaPipe et la géométrie des repères derrière des fonctions \code{is\_*()} lisibles. Le développeur d'application lit un état et écrit sa logique.
  \item \textbf{Développement piloté par specs.} Chaque application a ses exigences au format EARS, un design avec des propriétés de correction reliées aux exigences (« \emph{For any} \dots{} \textbf{Validates: Requirements 1.4} ») et un plan de tâches relié aux exigences.
  \item \textbf{Steering prescriptif.} Les fichiers imposent la structure des dossiers, le nommage, l'architecture (« Direct Controller ») et un modèle de code. L'agent produit ainsi du code homogène.
  \item \textbf{Hooks d'intégration.} L'extension du cœur (un nouveau geste) déclenche automatiquement toutes les modifications d'intégration. C'est leur usage de Kiro le plus original.
  \item \textbf{Présentation soignée.} Bannière, GIF de démonstration, instructions courtes, deux applications au thème Halloween.
\end{enumerate}

\subsection{Limites observées}

\begin{itemize}[leftmargin=1.2em]
  \item \textbf{Frontière faible.} Le « contrat » entre cœur et applications se réduit à un dataclass de booléens. Rien ne vérifie qu'une application respecte le cœur, et le cœur ne propose aucun protocole.
  \item \textbf{Applications monolithiques.} Toute la logique d'une application tient dans un fichier, ce qui serait difficile à faire évoluer à plusieurs.
  \item \textbf{Tests non réalisés.} Les 36 tâches de tests optionnelles (29 tests de propriétés, 5 tests unitaires, 2 autres) sont restées décochées. Les propriétés sont écrites mais jamais vérifiées.
  \item \textbf{Hooks coûteux et fragiles.} Les trois hooks sont des prompts d'agent (consomment des crédits, résultat non déterministe) et utilisent l'ancien format.
\end{itemize}

\subsection{Ce que nous en tirons pour Wysp}

\begin{xltabular}{\textwidth}{@{}L{4.2cm}L{4.6cm}Y@{}}
\toprule
\textbf{Chez InteractionKit} & \textbf{Ce que Wysp reprend} & \textbf{Ce que Wysp fait en plus} \\
\midrule
Cœur qui cache MediaPipe & Le cœur cache toute l'intelligence tutorielle : LLM, tableau, voix, pédagogie, permissions & Un contrat versionné, typé, exporté en JSON Schema et en types TypeScript \\
Deux applications de thème différent & Deux tuteurs dans \file{apps/blender-tutor} et \file{apps/web-tutor} & Deux \emph{niveaux d'intégration} différents (natif et interface), preuve vérifiable que le cœur n'a pas changé (contrôle de frontière) \\
Specs EARS et propriétés & Même format, même traçabilité exigence $\rightarrow$ propriété $\rightarrow$ tâche & Les tâches de propriétés sont \textbf{obligatoires} et tournent en CI \\
Hooks d'intégration de gestes & Hook « nouvelle action d'adaptateur » qui câble registre, undo, fixture de contrat et README & Hooks déterministes (commandes Python) pour les contrôles fréquents ; prompts d'agent seulement pour la revue de sécurité \\
Steering prescriptif & Produit, technique, structure, tests & Règles de confiance et sécurité, guide du contrat, pédagogie, charte du panneau \\
MCP Context7 & Context7 pour \code{bpy} et Chrome & Playwright MCP, puis un pont MCP en lecture seule vers un tuteur en cours d'exécution \\
README visuel et onboarding & GIF, démarrage rapide, licence MIT & « Créez votre tuteur en quatre étapes », kit de contrat pour valider un nouvel adaptateur \\
\bottomrule
\end{xltabular}

\begin{keybox}[Leçon principale]
Le jury de la catégorie récompense un squelette qui supprime une complexité répétitive et dont la réutilisation est démontrée. Pour Wysp, la meilleure preuve est double : le contrôle de frontière qui montre qu'aucun code d'application n'entre dans le cœur, et le même kit de tests de contrat qui passe contre l'adaptateur factice, Blender et Chrome.
\end{keybox}

\section{Architecture technique}

\subsection{Vue d'ensemble}

Wysp tourne entièrement en local. Un processus Python, \textbf{Wysp Core}, sert le \textbf{panneau} (application React ouverte dans une fenêtre Chrome à côté du logiciel) et accepte les connexions des \textbf{adaptateurs}. Chaque adaptateur vit dans son application hôte : une extension Blender, une extension Chrome. Le cœur appelle un fournisseur LLM et un service de synthèse vocale configurés par variables d'environnement.

\begin{figure}[h]
\centering
\begin{tikzpicture}[
  font=\sffamily\scriptsize,
  box/.style={draw=#1, fill=#1!8, rounded corners=2pt, align=center, minimum height=0.62cm, inner sep=3pt},
  mod/.style={draw=pa!70, fill=white, rounded corners=1.5pt, align=center, minimum width=2.05cm, minimum height=0.55cm, inner sep=2pt},
  lab/.style={font=\sffamily\scriptsize\itshape, text=muted},
  arr/.style={-{Stealth[length=2.2mm]}, thick, color=surface},
  biarr/.style={{Stealth[length=2.2mm]}-{Stealth[length=2.2mm]}, thick, color=surface}]

  % Core
  \node[draw=pa, fill=pa!5, rounded corners=4pt, minimum width=9.0cm, minimum height=4.5cm] (core) at (0,0.5) {};
  \node[font=\sffamily\bfseries\small, text=pa] at (0,2.47) {Wysp Core (Python, FastAPI, asyncio)};
  \node[draw=pa!70, fill=pa!12, rounded corners=1.5pt, minimum width=8.7cm, minimum height=0.42cm, inner sep=1pt] (ctr) at (0,1.98) {contract : modèles Pydantic, prédicats, empreinte d'état (partagés avec les SDK et le kit)};
  \node[mod, minimum width=1.5cm, minimum height=2.85cm] (gw) at (-3.65,0.12) {gateway\\[2pt]\tiny JSON-RPC\\\tiny registre\\\tiny jetons};
  \node[mod, minimum width=1.5cm, minimum height=2.85cm] (ui) at (3.65,0.12) {ui\_api\\[2pt]\tiny /ws/ui\\\tiny /api/tts\\\tiny panneau};
  \node[mod, minimum width=1.7cm] (sess)  at (-1.9, 1.12) {session\\\tiny bus d'événements};
  \node[mod, minimum width=1.7cm] (orch)  at ( 0.0, 1.12) {orchestrator\\\tiny tours, outils};
  \node[mod, minimum width=1.7cm] (board) at ( 1.9, 1.12) {board\\\tiny projection};
  \node[mod, minimum width=1.7cm] (trust) at (-1.9, 0.12) {trust\\\tiny porte, undo, audit};
  \node[mod, minimum width=1.7cm] (verif) at ( 0.0, 0.12) {verification\\\tiny prédicats, preuves};
  \node[mod, minimum width=1.7cm] (narr)  at ( 1.9, 0.12) {narration\\\tiny segments};
  \node[mod, minimum width=1.7cm] (ped)   at (-1.9,-0.88) {pedagogy\\\tiny leçons, apprenant};
  \node[mod, minimum width=1.7cm] (llm)   at ( 0.0,-0.88) {llm\\\tiny fournisseurs};
  \node[mod, minimum width=1.7cm] (store) at ( 1.9,-0.88) {store\\\tiny \textasciitilde/.wysp};

  % Adapters
  \node[box=pc, minimum width=2.6cm, minimum height=1.35cm] (bl) at (-6.9, 0.95) {};
  \node[font=\sffamily\bfseries\scriptsize, text=pc] at (-6.9, 1.42) {Blender 5.2 LTS};
  \node[box=pc, fill=white, minimum width=2.3cm] (bla) at (-6.9, 0.8) {Extension Wysp\\\tiny bpy, timers, gpu};
  \node[box=pc, minimum width=2.6cm, minimum height=1.35cm] (ch) at (-6.9,-0.75) {};
  \node[font=\sffamily\bfseries\scriptsize, text=pc] at (-6.9,-0.28) {Chrome};
  \node[box=pc, fill=white, minimum width=2.3cm] (cha) at (-6.9,-0.9) {Extension MV3\\\tiny service worker, script};

  % Panel
  \node[box=pb, minimum width=2.6cm, minimum height=2.85cm] (panel) at (6.8, 0.12) {};
  \node[font=\sffamily\bfseries\scriptsize, text=pb] at (6.8, 1.3) {Panneau (React)};
  \node[box=pb, fill=white, minimum width=2.3cm] at (6.8, 0.75) {Tableau\\\tiny React Flow, ELK};
  \node[box=pb, fill=white, minimum width=2.3cm] at (6.8, 0.0) {Sous-titres, voix};
  \node[box=pb, fill=white, minimum width=2.3cm] at (6.8,-0.75) {Permissions\\\tiny exercices, leçons};

  % External
  \node[box=okgreen, minimum width=2.4cm] (prof) at (-1.9,-2.75) {Profils YAML\\\tiny apps/*/profile};
  \node[box=muted, minimum width=2.4cm] (prov) at (0.75,-2.75) {Fournisseur LLM\\\tiny appels d'outils};
  \node[box=muted, minimum width=2.4cm] (tts) at (3.65,-2.75) {Synthèse vocale\\\tiny clé côté serveur};

  % Arrows
  \draw[biarr] (bla.east) -- (gw.west |- bla.east);
  \draw[biarr] (cha.east) -- (gw.west |- cha.east);
  \node[font=\sffamily\tiny, text=muted, align=center] at (-5.05,-0.05) {JSON-RPC\\WebSocket};
  \draw[biarr] (ui.east |- 0,0.75) -- node[lab, above] {/ws/ui} (panel.west |- 0,0.75);
  \draw[arr] (prof.north) -- (ped.south);
  \draw[arr] (llm.south) -- (llm.south |- prov.north);
  \draw[arr] (ui.south) -- (tts.north);
  \node[lab] at (0.42,-1.78) {HTTPS};
  \node[lab] at (4.12,-1.78) {HTTPS};
\end{tikzpicture}
\caption{Vue d'ensemble. Seuls les adaptateurs (orange) et les profils (vert) dépendent de l'application.}
\end{figure}

\begin{xltabular}{\textwidth}{@{}L{3.1cm}L{2.2cm}Y L{1.2cm}@{}}
\toprule
\textbf{Composant} & \textbf{Langage} & \textbf{Rôle} & \textbf{Resp.} \\
\midrule
Contrat & Python (Pydantic) & Modèles, export JSON Schema, prédicats, empreinte d'état & \PA \\
Wysp Core & Python 3.13 & Passerelle, session, orchestrateur, LLM, tableau, narration, confiance, pédagogie, vérification & \PA \\
Kit de contrat, adaptateur factice & Python & Tests contre chaque adaptateur, application Kanban fictive & \PA \\
Panneau & TypeScript, React 19 & Tableau, sous-titres, voix, dialogues & \PB \\
Tableau de bord de démo & TypeScript, React & Application web cible « Haunted Shop Analytics » & \PB \\
SDK adaptateur & Python et TypeScript & Client WebSocket, registre d'actions, validation & \PC \\
Tuteur Blender & Python (bpy) & Extension Blender 5.2 et scène de démo (\PC), contenu du profil (\PB) & \PC{} \PB \\
Tuteur web & TypeScript & Extension Chrome MV3 (\PC), profil web (\PB) & \PC{} \PB \\
\bottomrule
\end{xltabular}

\subsection{Principes directeurs}

\begin{enumerate}[leftmargin=1.2em]
  \item \textbf{Le cœur ignore les applications.} Aucun import d'hôte (\code{bpy}, API du navigateur) ni vocabulaire d'application dans \file{core/}. Un script vérifie cette règle à chaque modification, en CI et via un hook Kiro.
  \item \textbf{Le LLM décide quoi dire et quoi montrer ; le cœur décide ce qui est vrai et ce qui est permis.} Le LLM n'a aucun outil d'exécution directe. Il propose ; la porte de permission décide ; l'adaptateur exécute ; le vérificateur contrôle.
  \item \textbf{Une seule source de vérité pour le contrat.} Modèles Pydantic, puis JSON Schema, puis types TypeScript générés. La CI échoue si les fichiers générés ne correspondent plus.
  \item \textbf{Les connaissances d'application sont des données.} Le profil (YAML) décrit compétences, concepts, leçons et exercices. Les prompts du cœur restent génériques.
  \item \textbf{Tout est testable sans l'application réelle} : adaptateur factice, LLM scénarisé, enregistrements d'échanges.
  \item \textbf{Local d'abord.} Écoute sur \code{127.0.0.1}, jeton d'appairage, liste blanche d'origines, clés côté serveur.
\end{enumerate}

\subsection{Le contrat d'adaptateur}

\subsubsection{Transport et session}

Les adaptateurs sont \textbf{clients} et le cœur est \textbf{serveur}. Le Python embarqué de Blender et le service worker d'une extension Chrome savent ouvrir une connexion WebSocket sortante, alors qu'héberger un serveur y serait difficile. Les messages sont du JSON-RPC 2.0 dans les deux sens sur la même connexion.

\begin{enumerate}[leftmargin=1.2em]
  \item L'adaptateur envoie \code{wysp.hello \{manifest, token\}}. Le cœur vérifie le jeton (fichier \file{~/.wysp/token}), la version majeure du contrat et le manifeste.
  \item L'adaptateur envoie \code{wysp.heartbeat} au moins toutes les 20\,s. Cette fréquence maintient aussi en vie le service worker Chrome (Chrome 116 et suivants). Le cœur déclare l'adaptateur déconnecté après 45\,s de silence.
  \item Le cœur appelle les méthodes \code{adapter.*}. L'adaptateur envoie des notifications \code{event.*}.
\end{enumerate}

{\small\setlength{\tabcolsep}{4pt}
\begin{xltabular}{\textwidth}{@{}L{3.7cm}L{3.0cm}L{2.0cm}L{0.9cm}Y@{}}
\toprule
\textbf{Méthode} & \textbf{Résultat} & \textbf{Capacité} & \textbf{Budget} & \textbf{Correspondance avec la note de concept} \\
\midrule
\code{adapter.observe} & Observation (graphe d'entités) & \code{observe} & 1\,s & \code{observe()} \\
\code{adapter.get\_context} & Contexte compact ($\leq$ 8\,Ko) & \code{context} & 0,5\,s & \code{get\_context()} \\
\code{adapter.get\_actions} & Catalogue d'actions & \code{actions} & 0,5\,s & \code{get\_actions()} \\
\code{adapter.execute} & Résultat, jeton d'annulation & \code{execute} & 5\,s & \code{execute(action)} \\
\code{adapter.undo} & Résultat & \code{undo} & 5\,s & ajouté : réversibilité \\
\code{adapter.highlight} & accusé & \code{highlight} & 0,5\,s & ajouté : montrer dans l'application \\
\code{adapter.verify} & Preuves & \code{verify.native} & 2\,s & \code{verify(goal)}, optionnel \\
\code{event.state\_changed} & notification & \code{events.state} & -- & événements \\
\code{event.selection\_changed}\newline \code{event.user\_action}\newline \code{event.navigation} & notifications & \code{events.*} & -- & événements \\
\bottomrule
\end{xltabular}}

\code{verify(goal)} et \code{evaluate(exercise)} de la note de concept sont surtout réalisés \textbf{par le cœur} : il évalue des prédicats sur une observation fraîche. Un adaptateur peut fournir \code{adapter.verify} quand une vérification native est plus fiable (par exemple un calcul de géométrie dans Blender). Ce choix évite de réécrire la logique de vérification dans chaque adaptateur.

\subsubsection{Le modèle d'observation : un graphe d'entités}

Une observation contient des \textbf{entités} identifiées par une référence stable de type URI, avec un \code{kind} déclaré dans le manifeste, un parent, des propriétés JSON et une confiance entre 0 et 1. Elle contient aussi des \textbf{liens} typés et un \textbf{focus} (sélection ou élément actif).

\begin{lstlisting}[language={}]
{"observation_id": "o-118", "state_version": 42, "scope": {"level": "focus"},
 "entities": [
  {"ref": "blender://object/Pumpkin", "kind": "object", "label": "Pumpkin",
   "props": {"type": "MESH", "modifiers_count": 2}},
  {"ref": "blender://object/Pumpkin/modifier/Subdivision", "kind": "modifier",
   "label": "Subdivision", "parent": "blender://object/Pumpkin",
   "props": {"type": "SUBSURF", "index": 0, "levels": 2, "show_viewport": true}},
  {"ref": "blender://object/Pumpkin/modifier/Solidify", "kind": "modifier",
   "label": "Solidify", "parent": "blender://object/Pumpkin",
   "props": {"type": "SOLIDIFY", "index": 1, "thickness": 0.02}}],
 "links": [], "focus": ["blender://object/Pumpkin"], "truncated": false}
\end{lstlisting}

La même structure décrit une page web : \code{web://…/region\textasciitilde{}filters/combobox\textasciitilde{}region} de type \code{control}, avec \code{role}, \code{name}, \code{value} et \code{states}. Le cœur ne connaît ni les objets 3D ni les champs de formulaire : il manipule des entités. Les limites sont 500 entités et 256\,Ko par observation, avec un indicateur \code{truncated} et des comptes par type.

\subsubsection{Actions, risques et post-conditions}

Chaque adaptateur publie un catalogue \textbf{fermé} d'actions. Une action déclare un identifiant, une description destinée au LLM, un schéma JSON de paramètres, un niveau de risque, des post-conditions et, si l'utilisateur peut la faire lui-même, un chemin d'interface (\code{guide.ui\_path}) utilisé en mode guidé. Aucune action ne peut évaluer du code arbitraire.

\begin{xltabular}{\textwidth}{@{}L{2.3cm}Y L{4.3cm}@{}}
\toprule
\textbf{Risque} & \textbf{Politique par défaut} & \textbf{Exemples} \\
\midrule
\code{read\_only} & Exécution automatique, journalisée & lecture de propriétés \\
\code{ui\_state} & Exécution automatique, journalisée & sélectionner un objet, changer d'onglet, faire défiler \\
\code{reversible} & Approbation : clic, « oui » vocal ou autorisation limitée à la leçon. Jeton d'annulation obligatoire & ajouter un modificateur, remplir un champ, choisir un filtre \\
\code{irreversible} & Clic explicite avec aperçu, jamais par la voix, interdit au niveau visuel & appliquer un modificateur, cliquer un bouton quelconque \\
\bottomrule
\end{xltabular}

\subsubsection{Langage de prédicats}

Les post-conditions, les critères de réussite des exercices et les objectifs de vérification utilisent un seul langage, évalué par le cœur : des comparaisons sur une expression JMESPath, combinées par \code{all}, \code{any} et \code{not}.

\begin{lstlisting}[language={}]
{"all": [
  {"path": "entities[?kind=='modifier' && parent=='${bindings.target}' && props.type=='BEVEL'] | length(@)",
   "op": "gte", "value": 1},
  {"path": "entities[?kind=='modifier' && parent=='${bindings.target}' && props.type=='BEVEL'].props.segments | [0]",
   "op": "eq", "value": 3}]}
\end{lstlisting}

Les paramètres ne sont substitués que dans des positions autorisées, avec échappement, pour empêcher une valeur d'injecter de la syntaxe JMESPath (propriété CON-P11). L'évaluateur renvoie la valeur réelle et la valeur attendue : ce sont les éléments de la \textbf{preuve} affichée à l'utilisateur.

\subsubsection{Empreinte d'état, erreurs et versions}

\begin{itemize}[leftmargin=1.2em]
  \item \textbf{Empreinte d'état} : SHA-256 d'un JSON canonique des entités (sans propriétés volatiles, flottants arrondis) et des liens triés. Elle définit « l'état est restauré » pour l'annulation et « rien n'a changé » pour les appels invalides.
  \item \textbf{Erreurs} : codes JSON-RPC standard et codes Wysp de \code{-32010} à \code{-32019} (non supporté, action inconnue, cible introuvable, précondition, échec, annulation impossible, délai, occupé, version incompatible, non autorisé).
  \item \textbf{Versions} : \code{contract\_version} commence à 1.0.0 au jalon J0. Ajout optionnel : version mineure. Changement cassant : version majeure, accord des trois personnes.
\end{itemize}

\subsubsection{Kit de tests de contrat}

Le kit (\code{wysp-contract --tutor <nom>}) joue le rôle du cœur : il démarre un serveur de test, lance l'adaptateur selon \file{apps/<tuteur>/contract.fixture.yaml}, puis vérifie la poignée de main, le manifeste, les invariants des observations, le catalogue, les appels sûrs (exécution, post-conditions, annulation, empreinte), les appels invalides, l'étiquetage d'origine des événements et les budgets de temps. Il enregistre les échanges. En mode \code{--offline}, il revalide ces enregistrements contre les schémas courants : chacun sait ainsi quel tuteur casserait un changement de contrat, sans lancer Blender ni Chrome.

\subsection{Wysp Core}

\subsubsection{Modules}

\begin{xltabular}{\textwidth}{@{}L{2.7cm}Y@{}}
\toprule
\textbf{Module} & \textbf{Responsabilité} \\
\midrule
\file{contract/} & Modèles Pydantic du contrat et du protocole du panneau, export des schémas, prédicats, empreinte d'état. Seul module importé par les SDK et le kit. \\
\file{gateway/} & Serveur JSON-RPC \code{/ws/adapter}, \code{AdapterProxy} typé avec budgets de temps, registre des adaptateurs, surveillance des battements, contrôle d'origine et de jeton. \\
\file{session/} & Session de tutorat (adaptateur actif, profil, tableau, narration, propositions), bus d'événements asyncio, mémoire de conversation avec résumé glissant, constructeur de contexte. \\
\file{orchestrator/} & Boucle de tour (\code{TurnRunner}), outils offerts au LLM, prompts génériques, annulation. \\
\file{llm/} & Interface \code{LLMProvider}, \code{ScriptedLLM} (scénarios rejouables), fournisseur compatible OpenAI (OpenAI, Groq, Ollama local…), Bedrock en option. \\
\file{board/} & Modèle visuel, opérations et correctifs versionnés, projection des entités, validateur de provenance, regroupement au-delà de 60 nœuds. \\
\file{narration/} & Segments, repères de synchronisation, curseur de lecture, interruption et reprise. \\
\file{trust/} & Politique de risque, porte de permission, autorisations de leçon, pile d'annulation, journal d'audit. \\
\file{verification/} & Observation fraîche, évaluation des prédicats, preuves. \\
\file{pedagogy/} & Profils, graphe de compétences, leçons, exercices, modèle de l'apprenant, recommandations. \\
\file{ui_api/}, \file{store/} & Protocole du panneau, proxy de synthèse vocale, page d'appairage ; persistance JSON/JSONL dans \file{~/.wysp/} avec masquage des secrets. \\
\bottomrule
\end{xltabular}

\subsubsection{La boucle de tour et les outils du LLM}

Un \textbf{tour} traite un message de l'utilisateur. Le cœur construit un contexte d'au plus 6\,000 jetons : prompt système générique, manifeste (types d'entités, catalogue d'actions), amorce du profil, contexte de l'adaptateur et observation du focus entourés de balises \code{<app\_data>} (contenu traité comme des données, jamais comme des instructions), état de l'apprenant, leçon en cours et conversation récente. Le LLM répond uniquement par des appels d'outils, au plus 8 tours d'outils par message.

\begin{xltabular}{\textwidth}{@{}L{3.0cm}Y L{2.3cm}@{}}
\toprule
\textbf{Outil} & \textbf{Effet} & \textbf{Capacité requise} \\
\midrule
\code{observe} & Résumé compact d'une observation & \code{observe} \\
\code{show\_entities} & Projection d'entités réelles sur le tableau (nœuds \pobs) & \code{observe} \\
\code{add\_concept}, \code{annotate} & Nœuds et annotations \pinf{} ou \pill, jamais \pobs & -- \\
\code{narrate} & Segments de narration avec repères (focus, surbrillance, étape) & -- \\
\code{highlight} & Mise en évidence dans l'application & \code{highlight} \\
\code{propose\_action} & Proposition à la porte de permission, mode \code{guide} ou \code{do} & \code{actions}, \code{execute} \\
\code{start\_exercise} & Démarrage d'un exercice du profil si ses liaisons se résolvent & \code{observe} \\
\code{lookup\_concept} & Entrée du glossaire du profil & -- \\
\code{ask\_clarification} & Question avec au plus 4 réponses proposées & -- \\
\bottomrule
\end{xltabular}

Les outils dont la capacité manque au manifeste ne sont pas proposés (propriété CORE-P10). Il n'existe aucun outil « exécuter » : c'est la garantie centrale du modèle de confiance.

\subsection{Modèle visuel et tableau}

Le tableau est un \textbf{modèle structuré}, pas une image : nœuds (avec \code{ref}, \code{kind}, provenance, confiance, propriétés à afficher), arêtes, étapes, annotations, focus et disposition. Le cœur publie des correctifs versionnés (\code{board.patch \{base\_version, version, ops\}}) ; le panneau les applique dans l'ordre et demande un instantané en cas de trou. La même fonction d'application existe en Python et en TypeScript, vérifiée par des fixtures communes (CORE-P3, PAN-P1).

\begin{itemize}[leftmargin=1.2em]
  \item \textbf{Projection.} Une entité devient un nœud d'identifiant \code{n:<ref>} ; les relations parent deviennent des arêtes \code{contains} ; les liens d'observation deviennent des arêtes typées. Le profil choisit les propriétés affichées par type (\code{board.props}). La projection est déterministe (CORE-P12), ce qui stabilise la disposition.
  \item \textbf{Validateur de provenance.} Tout nœud marqué \pobs{} dont la référence n'apparaît dans aucune observation de la session est rétrogradé en \pinf{} et journalisé (CORE-P1).
  \item \textbf{Dispositions.} Quatre suffisent pour les deux tuteurs et couvrent les exemples de la note de concept.
\end{itemize}

\begin{xltabular}{\textwidth}{@{}L{2.0cm}Y Y@{}}
\toprule
\textbf{Disposition} & \textbf{Blender} & \textbf{Web} \\
\midrule
\code{stack} & pile de modificateurs, ordre d'évaluation & ordre des filtres appliqués \\
\code{tree} & hiérarchie de la scène (collections, objets) & régions de la page, contrôles d'un formulaire \\
\code{graph} & objets, matériaux et références & dépendances entre filtres, KPI et tableau \\
\code{sequence} & étapes guidées d'une leçon & étapes de création d'un rapport \\
\bottomrule
\end{xltabular}

\begin{xltabular}{\textwidth}{@{}L{2.4cm}Y L{4.4cm}@{}}
\toprule
\textbf{Provenance} & \textbf{Style du nœud} & \textbf{Marqueur dans les sous-titres} \\
\midrule
\pobs & bordure pleine 2\,px, opacité complète, icône « œil » & aucun \\
\pinf & bordure pointillée, opacité 85\,\%, badge de confiance & « probablement » \\
\pill & fond hachuré, bordure grise, étiquette « exemple » & « par exemple » \\
\bottomrule
\end{xltabular}

\subsection{Voix et narration}

Nous retenons une chaîne en cascade plutôt qu'un modèle vocal de bout en bout : transcription dans le navigateur, raisonnement en texte dans le cœur, synthèse vocale par segments. Le texte reste ainsi la source de vérité : sous-titres, provenance et repères de synchronisation sont exacts, et la voix peut tomber en panne sans casser le tutorat.

\begin{figure}[h]
\centering
\begin{tikzpicture}[font=\sffamily\scriptsize, node distance=0.55cm,
  st/.style={draw=#1, fill=#1!8, rounded corners=2pt, align=center, minimum height=0.8cm, minimum width=1.9cm},
  arr/.style={-{Stealth[length=2mm]}, thick, color=surface}]
  \node[st=pb] (mic) {Micro\\\tiny appui pour parler};
  \node[st=pb, right=of mic] (stt) {Transcription\\\tiny Web Speech API};
  \node[st=pa, right=of stt] (core) {Cœur : tour\\\tiny LLM + outils};
  \node[st=pa, right=of core] (seg) {Segments\\\tiny $\leq$ 300 car. + repères};
  \node[st=pb, right=of seg] (tts) {Synthèse\\\tiny /api/tts ou navigateur};
  \node[st=pb, right=of tts] (spk) {Audio et\\sous-titres};
  \draw[arr] (mic) -- (stt); \draw[arr] (stt) -- (core); \draw[arr] (core) -- (seg); \draw[arr] (seg) -- (tts); \draw[arr] (tts) -- (spk);
  \draw[arr, color=danger, dashed] (spk.south) -- ++(0,-0.55) -| node[pos=0.25, below, text=danger] {interruption : parole détectée ou Échap $\rightarrow$ \code{user.interrupt}, arrêt en moins de 150\,ms} (mic.south);
\end{tikzpicture}
\caption{Chaîne vocale. En violet le panneau, en bleu le cœur.}
\end{figure}

\begin{itemize}[leftmargin=1.2em]
  \item \textbf{Appui pour parler} (barre d'espace ou bouton) par défaut : fiable en démonstration, pas de déclenchement intempestif.
  \item \textbf{Interruption} : pendant la narration, une détection d'énergie vocale (WebAudio, seuil calibré, annulation d'écho) ou la touche Échap arrête l'audio et envoie \code{user.interrupt}. Le cœur annule le tour et mémorise le dernier segment joué. « Continue » rejoue exactement les segments non entendus (CORE-P7, CORE-P8).
  \item \textbf{Clés} : la synthèse passe par \code{/api/tts} pour que la clé reste côté serveur ; la synthèse du navigateur sert de repli.
  \item \textbf{Latence visée} : premier son moins de 2,5\,s après la fin de la question, grâce au découpage en segments et à la prélecture du segment suivant.
\end{itemize}

\subsection{Pédagogie et vérification}

Le moteur pédagogique lit le \textbf{profil} de l'application : compétences avec prérequis, glossaire de concepts reliés à la documentation officielle, leçons en étapes (\code{explain}, \code{show}, \code{guide}, \code{do}, \code{check}) et exercices.

\begin{lstlisting}[language={}]
exercises:
  - id: bevel_segments
    prompt: {en: "Give the bevel 3 segments so the edge looks round."}
    skills: [bevel_modifier]
    bindings: {target: "$focus.object"}
    success: {all: [ ...predicates on the observation... ]}
    hints: [{en: "Look at the Modifier Properties tab (blue wrench)."},
            {en: "Find the Segments field in the Bevel panel."}]
    max_attempts: 3
\end{lstlisting}

\begin{itemize}[leftmargin=1.2em]
  \item \textbf{Exercices contextualisés} : un exercice n'est proposé que si ses liaisons se résolvent sur la scène ou la page réelle et si son critère n'est pas déjà satisfait (PED-P5, PED-P6).
  \item \textbf{Vérification} : à chaque changement d'origine utilisateur, le moteur évalue les prédicats sur une observation fraîche et enregistre une preuve (prédicat, valeur réelle, valeur attendue, identifiant d'observation, horodatage).
  \item \textbf{Trois niveaux de preuve} : \code{verified} (constaté dans l'application, poids 1,0), \code{answered} (réponse à un quiz, 0,6), \code{declared} (« c'est fait » non vérifiable, 0,3). Le panneau affiche toujours la différence. Les actions faites par Wysp ne comptent jamais (PED-P4).
  \item \textbf{Modèle de l'apprenant} : traçage bayésien des connaissances simplifié (glissement 0,1, chance 0,2, apprentissage 0,15), pondéré par le type de preuve et le nombre d'indices. Une compétence est acquise à 0,85. La profondeur des explications suit la maîtrise : bases sous 0,4, pratique guidée jusqu'à 0,85, défi au-delà.
\end{itemize}

\subsection{Confiance et sécurité}

\begin{figure}[h]
\centering
\begin{tikzpicture}[font=\sffamily\scriptsize,
  s/.style={draw=#1, fill=#1!8, rounded corners=2pt, align=center, minimum height=0.66cm, minimum width=1.75cm},
  arr/.style={-{Stealth[length=2mm]}, thick, color=surface},
  lbl/.style={font=\sffamily\tiny, text=muted, fill=white, inner sep=1pt}]
  \node[s=pa] (p) at (0,0) {proposée\\\tiny paramètres validés};
  \node[s=pa] (pol) at (2.4,0) {politique\\\tiny selon le risque};
  \node[s=okgreen] (auto) at (5.0,0.95) {approuvée\\\tiny auto ou autorisation};
  \node[s=pumpkin] (wait) at (5.0,-0.95) {en attente\\\tiny délai 60\,s};
  \node[s=muted] (deny) at (5.0,-2.35) {refusée\\\tiny ou expirée};
  \node[s=okgreen] (ex) at (7.7,0) {exécution\\\tiny adaptateur};
  \node[s=pa] (ver) at (10.1,0) {vérification\\\tiny post-conditions};
  \node[s=okgreen] (done) at (12.6,0.95) {réussie\\\tiny preuve enregistrée};
  \node[s=danger] (undo) at (12.6,-0.95) {échec\\\tiny annulation auto};
  \draw[arr] (p) -- (pol);
  \draw[arr] (pol) -- (auto);
  \draw[arr] (pol) -- (wait);
  \draw[arr] (auto) -- (ex);
  \draw[arr] (wait) -- node[lbl, pos=0.45, below right=-1pt] {clic ou « oui »} (ex);
  \draw[arr] (wait) -- node[lbl, right] {refus} (deny);
  \draw[arr] (ex) -- (ver);
  \draw[arr] (ver) -- (done);
  \draw[arr] (ver) -- (undo);
\end{tikzpicture}
\caption{Cycle de vie d'une proposition d'action.}
\end{figure}

\begin{itemize}[leftmargin=1.2em]
  \item \textbf{Provenance} : seule la projection d'observations crée des nœuds \pobs ; le validateur rattrape le reste. Au niveau d'intégration visuel, la confiance des entités est inférieure à 1 et affichée.
  \item \textbf{Porte de permission} : politique de la table des risques, autorisations limitées à une leçon pour les actions réversibles, clic obligatoire et aperçu pour les actions irréversibles, expiration à 60\,s. Une surcharge de politique par profil ne peut être que plus stricte (CORE-P4).
  \item \textbf{Vérifier ou annuler} : après exécution, nouvelle observation et évaluation des post-conditions ; en cas d'échec, annulation automatique et rapport avec preuve (CORE-P6). Chaque action Blender pousse aussi une étape d'annulation native, donc Ctrl+Z fonctionne.
  \item \textbf{Sécurité locale} : écoute sur \code{127.0.0.1} uniquement, contrôle de l'en-tête \code{Origin}, jeton d'appairage, aucune clé dans le navigateur, journaux masqués (CORE-P11).
  \item \textbf{Vie privée (web)} : extension active seulement sur les origines autorisées ; jamais de mot de passe, carte bancaire, code à usage unique, champ caché, cookie ou stockage dans les observations (WEB-P1, WEB-P6).
  \item \textbf{Injection de prompt} : le texte d'une page ou d'une scène est une donnée. Une page piégée peut au pire faire \emph{proposer} une action ; la porte et le catalogue fermé empêchent qu'elle s'exécute sans l'utilisateur.
  \item \textbf{Audit} : propositions, décisions (qui, par quel canal), exécutions, vérifications et annulations dans \file{~/.wysp/audit/*.jsonl} (CORE-P5).
\end{itemize}

\subsection{Adaptateur Blender (intégration native)}

L'adaptateur est une extension Blender 5.2 LTS (Python 3.13, \file{blender_manifest.toml}, permissions \code{network} et \code{files} justifiées, dépendances pures Python embarquées en wheels). Licence GPL-3.0-or-later pour ce dossier, comme l'exige l'écosystème des extensions Blender ; le reste du dépôt est sous MIT.

\begin{figure}[h]
\centering
\begin{tikzpicture}[font=\sffamily\scriptsize,
  t/.style={draw=#1, fill=#1!6, rounded corners=3pt, minimum width=5.6cm, minimum height=2.5cm},
  s/.style={draw=#1, fill=white, rounded corners=2pt, align=center, minimum height=0.55cm, minimum width=2.2cm},
  arr/.style={-{Stealth[length=2mm]}, thick, color=surface}]
  \node[t=muted] at (0,0) {}; \node[font=\sffamily\bfseries\scriptsize] at (0,1.0) {Fil réseau (websocket-client)};
  \node[t=pc] at (7.2,0) {}; \node[font=\sffamily\bfseries\scriptsize, text=pc] at (7.2,1.0) {Fil principal de Blender};
  \node[s=muted] (rx) at (-1.3,0.25) {réception\\\tiny requête JSON-RPC};
  \node[s=muted] (tx) at (-1.3,-0.65) {envoi\\\tiny réponse, événement};
  \node[s=muted] (hb) at (1.35,-0.2) {battement\\\tiny 15\,s, sans bpy};
  \node[s=pc, minimum width=2.5cm] (pump) at (5.75,0.25) {pompe \texttt{\tiny bpy.app.timers}\\\tiny 50\,ms, budget 20\,ms};
  \node[s=pc, minimum width=2.0cm] (h) at (8.75,0.25) {gestionnaire\\\tiny utilise bpy};
  \node[s=pc, minimum width=2.6cm] (dg) at (7.2,-0.65) {\texttt{\tiny depsgraph\_update\_post}\\\tiny anti-rebond 250\,ms};
  \draw[arr] (rx.east) -- node[above, font=\sffamily\tiny] {file d'entrée} (pump.west);
  \draw[arr] (pump) -- (h);
  \draw[arr] (dg.west) -- node[below, font=\sffamily\tiny] {file de sortie} (tx.east);
\end{tikzpicture}
\caption{Modèle de threads : \code{bpy} n'est jamais appelé hors du fil principal (BLD-P1).}
\end{figure}

\begin{itemize}[leftmargin=1.2em]
  \item \textbf{Observation} : entités \code{scene}, \code{collection}, \code{object}, \code{modifier}, \code{material}. Pour les modificateurs, extraction \emph{générique} des propriétés RNA modifiables de types simples (\code{mod.bl\_rna.properties}) : aucun code par type de modificateur, et tous les types de Blender 5.2 sont couverts (BLD-P3). Focus : objet actif et sélection.
  \item \textbf{Actions} : neuf actions, une par fichier. API de données plutôt qu'opérateurs (\code{obj.modifiers.new}, \code{obj.modifiers.move}) ; \code{temp\_override} quand un opérateur est inévitable.
  \item \textbf{Annulation} : enregistrement compensatoire (opération inverse) pour chaque action réversible, plus \code{bpy.ops.ed.undo\_push} pour que Ctrl+Z fonctionne aussi (BLD-P2).
  \item \textbf{Événements} : gestionnaire \code{depsgraph\_update\_post} avec anti-rebond, comparaison d'empreintes par entité, événements dérivés (« modificateur ajouté »), \code{bpy.msgbus} pour l'objet actif, réabonnement au chargement d'un fichier. Une fenêtre d'étiquetage autour de chaque action Wysp donne \code{origin: wysp} (BLD-P4).
  \item \textbf{Mise en évidence} : contour pulsé de la boîte englobante dans la vue 3D (module \code{gpu}), bascule de l'éditeur de propriétés sur l'onglet Modificateurs.
  \item \textbf{Mode sans interface} : \file{demo/run_adapter.py} lance l'adaptateur avec \code{blender -b} et une pompe manuelle, pour le kit de contrat et la CI.
\end{itemize}

{\small
\begin{xltabular}{\textwidth}{@{}L{5.0cm}L{1.9cm}Y@{}}
\toprule
\textbf{Action} & \textbf{Risque} & \textbf{Chemin d'interface (mode guidé)} \\
\midrule
\code{select\_object}, \code{show\_properties\_tab} & ui\_state & clic dans l'Outliner ; icône d'onglet de l'éditeur de propriétés \\
\code{add\_modifier} & reversible & Propriétés › Modificateurs (clé) › Add Modifier › catégorie › type \\
\code{set\_modifier\_property} & reversible & panneau du modificateur › champ \\
\code{remove\_modifier}, \code{move\_modifier} & reversible & menu de l'en-tête du modificateur \\
\code{toggle\_modifier\_visibility} & reversible & icône écran de l'en-tête \\
\code{assign\_material} & reversible & Propriétés › Matériau › emplacement \\
\code{apply\_modifier} & irreversible & menu › Apply (Ctrl+A), avec aperçu du nombre de sommets \\
\bottomrule
\end{xltabular}}

\subsection{Adaptateur web (intégration par l'interface)}

\begin{figure}[h]
\centering
\begin{tikzpicture}[font=\sffamily\scriptsize,
  s/.style={draw=#1, fill=#1!8, rounded corners=2pt, align=center, minimum height=0.9cm, minimum width=2.6cm},
  arr/.style={{Stealth[length=2mm]}-{Stealth[length=2mm]}, thick, color=surface}]
  \node[s=pa] (core) at (0,0) {Wysp Core\\\tiny /ws/adapter};
  \node[s=pc] (sw) at (4.2,0) {Service worker\\\tiny client SDK, battement 20\,s};
  \node[s=pc] (cs) at (8.4,0) {Script de contenu\\\tiny instantané, actions, overlay};
  \node[s=muted] (dom) at (12.4,0) {Page autorisée\\\tiny DOM, accessibilité};
  \draw[arr] (core) -- node[above, font=\sffamily\tiny] {JSON-RPC} (sw);
  \draw[arr] (sw) -- node[above, font=\sffamily\tiny] {runtime.Port} (cs);
  \draw[arr] (cs) -- (dom);
\end{tikzpicture}
\caption{Extension Chrome MV3 : le service worker tient la connexion, le script de contenu fait tout le travail sur le DOM.}
\end{figure}

\begin{itemize}[leftmargin=1.2em]
  \item \textbf{Instantané orienté accessibilité} : rôles ARIA et HTML, noms et descriptions accessibles calculés avec \code{dom-accessibility-api}. Types d'entités : \code{page}, \code{region}, \code{heading}, \code{form}, \code{control}, \code{dialog}, \code{table}, \code{list}, \code{chart}, \code{status}. Options des listes (20 au plus), en-têtes et 10 premières lignes des tableaux, libellé des graphiques.
  \item \textbf{Localisateurs et références} : rôle + nom accessible dans la région nommée la plus proche, puis index, puis \code{data-testid}, puis chemin CSS. La référence encode le localisateur, donc elle reste stable tant que l'élément garde son rôle, son nom et sa région (WEB-P2, WEB-P3). Une cible ambiguë renvoie \code{TARGET\_NOT\_FOUND} avec des candidats.
  \item \textbf{Actions} : \code{scroll\_into\_view}, \code{focus\_element}, \code{follow\_link} (ui\_state, même origine seulement), \code{fill\_field}, \code{select\_option}, \code{set\_checked} (réversibles, l'annulation restaure la valeur précédente), \code{click\_element} (irréversible, aperçu de la cible). Les valeurs sont posées avec le \emph{setter} natif suivi des événements \code{input} et \code{change}, ce qui fonctionne avec React.
  \item \textbf{Origine des événements} : les événements synthétiques de Wysp ont \code{isTrusted === false}, ceux de l'utilisateur \code{true}. Les mutations du DOM sont attribuées à Wysp pendant la fenêtre d'exécution d'une action (WEB-P5).
  \item \textbf{Overlay} : racine Shadow DOM fermée, \code{pointer-events: none}, ignorée par l'instantané et l'observateur de mutations (WEB-P4).
  \item \textbf{Cible de démonstration} : « Haunted Shop Analytics » (\PB), tableau de bord accessible avec KPI, graphique, filtres (dates, région, catégories), tableau des commandes, dialogue « Create report » et liste des rapports, données déterministes et bouton de réinitialisation.
\end{itemize}

\subsection{Pile technique et structure du dépôt}

\begin{xltabular}{\textwidth}{@{}L{3.0cm}Y@{}}
\toprule
\textbf{Domaine} & \textbf{Choix} \\
\midrule
Cœur & Python 3.13, uv (espace de travail), FastAPI + uvicorn, Pydantic v2, jmespath, PyYAML, httpx ; pytest, Hypothesis, ruff \\
Panneau & React 19, Vite, \code{@xyflow/react} (React Flow 12), elkjs, Zustand ; Vitest, Testing Library, fast-check ; eslint, prettier \\
Blender & Blender 5.2 LTS, extension, websocket-client en wheel ; \code{bpy} de PyPI et fake-bpy-module pour les tests et l'éditeur \\
Navigateur & Chrome 116+, Manifest V3, dom-accessibility-api, Vite ; Playwright \\
Services & LLM avec appels d'outils (fournisseur compatible OpenAI par défaut, Bedrock en option), Web Speech API pour la transcription, synthèse vocale par proxy \\
CI & GitHub Actions : tests Python et TypeScript, export des schémas sans différence, contrôle de frontière, kit de contrat sur l'adaptateur factice \\
\bottomrule
\end{xltabular}

\begin{lstlisting}[language={}]
.kiro/            steering, specs, hooks, agents (commun, versionne)
core/             Wysp Core, paquet wysp_core                          (A)
schemas/          JSON Schemas generes                                 (A)
contract-tests/   kit de contrat, adaptateur factice Kanban, fixtures  (A)
sdk/python/       SDK adaptateur Python                                (C)
sdk/ts/           SDK TypeScript + types generes                       (C, types : A)
panel/            panneau React                                        (B)
apps/blender-tutor/   adapter/ (extension), profile/, demo/            (C)
apps/web-tutor/       extension/ (C), demo-dashboard/ (B), profile/ (B)
docs/             ce document, guide d'adaptateur, journal Kiro
scripts/kiro/     hook.py : controles appeles par les hooks Kiro et la CI
\end{lstlisting}

\subsection{Décisions d'architecture}

\begin{xltabular}{\textwidth}{@{}L{3.2cm}L{3.6cm}Y@{}}
\toprule
\textbf{Décision} & \textbf{Alternatives écartées} & \textbf{Raison} \\
\midrule
JSON-RPC 2.0 sur WebSocket, adaptateurs clients & HTTP REST ; gRPC ; MCP comme protocole d'adaptateur & Bidirectionnel (événements), simple dans Blender et dans un service worker, aucun serveur à héberger dans l'hôte. MCP reste utile côté développement (pont en lecture seule), mais ne porte ni événements ni permissions. \\
Pydantic source unique, JSON Schema, types TS générés & Écrire les schémas à la main ; TypeScript source & Une seule définition, validation à l'exécution des deux côtés, contrôle de dérive en CI. \\
Graphe d'entités générique & Un modèle par application & Le cœur, le tableau et les prédicats restent indépendants de l'application. \\
Vérification par prédicats évalués dans le cœur & \code{verify()} écrit dans chaque adaptateur & Moins de code par adaptateur, preuves uniformes ; \code{adapter.verify} reste possible. \\
Nœuds « observés » créés par projection & Laisser le LLM écrire le tableau & Garantit la règle de confiance de la note de concept. \\
Voix en cascade, appui pour parler & Modèle vocal temps réel de bout en bout & Sous-titres et provenance exacts, interruption maîtrisée, coût et dépendance moindres, repli possible. \\
Panneau web dans une fenêtre Chrome & Panneau natif dans Blender ; Electron & Une seule interface pour les deux tuteurs, c'est la preuve visuelle du cœur commun. \\
Catalogue d'actions fermé & Laisser le LLM écrire du \code{bpy} ou du JavaScript & Sécurité, vérifiabilité, annulation possible. \\
Stockage local en fichiers & Base de données & Suffisant pour un prototype, transparent pour l'utilisateur. \\
\bottomrule
\end{xltabular}


\section{Workflows}

Les quatre premiers workflows décrivent l'expérience de l'utilisateur ; ils sont identiques pour Blender et pour le web, seuls les contenus des observations et des actions changent. Les trois derniers décrivent notre façon de travailler.

% Notes pgf-umlsd : un libellé de \mess contenant des parenthèses doit être entre doubles accolades ;
% un tiret bas dans un libellé de \mess doit passer par \detokenize{...} ; pas de caractères actifs
% de babel français (: ; ! ? et guillemets) dans les libellés.

\subsection{Expliquer ce que l'utilisateur regarde}

Exemple : dans Blender, l'utilisateur sélectionne la citrouille et demande « Explique-moi ces modificateurs ».

\begin{center}
\resizebox{0.97\textwidth}{!}{%
\begin{sequencediagram}
\newthread{u}{Utilisateur}
\newinst[1.2]{p}{Panneau}
\newinst[1.6]{c}{Cœur}
\newinst[1.8]{l}{LLM}
\newinst[1.8]{a}{Adaptateur}
\mess{u}{{question (voix ou texte)}}{p}
\mess{p}{user.message}{c}
\begin{call}{c}{get\_context, observe(focus)}{a}{contexte, observation}
\end{call}
\begin{call}{c}{contexte + outils}{l}{show\_entities(refs, stack)}
\end{call}
\begin{callself}{c}{projection, validateur}{}
\end{callself}
\mess{c}{{board.patch (nœuds observed)}}{p}
\begin{call}{c}{résultat de l'outil}{l}{narrate(segments, repères)}
\end{call}
\mess{c}{narration.segment}{p}
\mess{c}{{adapter.highlight(ref)}}{a}
\mess{p}{voix, sous-titres, focus}{u}
\end{sequencediagram}}
\end{center}

Points clés : le LLM choisit \emph{quelles} entités montrer, la projection fabrique les nœuds à partir de l'observation réelle. Un concept général (par exemple « ordre d'évaluation ») est ajouté avec \code{add\_concept} en provenance \pinf{} ou \pill. Les repères des segments synchronisent la surbrillance sur le tableau et dans Blender avec la phrase prononcée.

\subsection{Action proposée, permission, exécution et vérification}

Exemple : « Arrondis les arêtes de la bougie. » Le LLM propose \code{add\_modifier(Candle, BEVEL)} en mode \code{do}.

\begin{center}
\resizebox{0.97\textwidth}{!}{%
\begin{sequencediagram}
\newthread{u}{Utilisateur}
\newinst[1.2]{p}{Panneau}
\newinst[1.6]{c}{Cœur et porte}
\newinst[1.8]{l}{LLM}
\newinst[1.8]{a}{Adaptateur}
\mess{l}{{\detokenize{propose_action(add_modifier, mode do)}}}{c}
\begin{callself}{c}{valider paramètres, politique}{risque reversible}
\end{callself}
\mess{c}{permission.request}{p}
\mess{p}{dialogue Approuver ou Refuser}{u}
\mess{u}{clic ou oui vocal}{p}
\mess{p}{permission.respond}{c}
\begin{call}{c}{adapter.execute}{a}{ok, undo\_token}
\end{call}
\mess{a}{{\detokenize{event.state_changed (origin wysp)}}}{c}
\begin{call}{c}{adapter.observe}{a}{observation fraîche}
\end{call}
\begin{callself}{c}{post-conditions}{preuve}
\end{callself}
\mess{c}{board.patch, evidence.added}{p}
\mess{c}{résultat vérifié}{l}
\end{sequencediagram}}
\end{center}

\begin{itemize}[leftmargin=1.2em]
  \item \textbf{Échec de vérification} : le cœur appelle \code{adapter.undo} avec le jeton, puis explique l'échec avec la valeur attendue et la valeur constatée.
  \item \textbf{Mode guidé} (préféré pour apprendre) : aucun \code{execute}. Le cœur publie les étapes de \code{guide.ui\_path} (« Propriétés › Modificateurs › Add Modifier › Bevel »), met en évidence chaque cible, et vérifie les post-conditions dès que l'adaptateur signale un changement d'origine utilisateur.
  \item \textbf{Annuler} : « annule ça » rejoue le dernier jeton de la pile d'annulation.
\end{itemize}

\subsection{Interruption et reprise}

\begin{enumerate}[leftmargin=1.2em]
  \item Pendant la narration, l'utilisateur parle (détection d'énergie) ou appuie sur Échap. Le panneau coupe l'audio et envoie \code{user.interrupt} (moins de 150\,ms).
  \item Le cœur annule le tour (jeton d'annulation vérifié entre chaque étape et pendant le flux du LLM), publie \code{narration.cancel} et garde le curseur : dernier segment entendu. Une action déjà commencée chez l'adaptateur n'est jamais coupée en deux ; sa vérification est enregistrée.
  \item Si l'utilisateur pose une question (« Pourquoi Solidify après Subdivision ? »), un nouveau tour démarre avec le nœud concerné dans le contexte. S'il dit « continue », le cœur republie exactement les segments non entendus, sans nouvel appel au LLM.
  \item Le bouton « Pourquoi ? » d'un nœud produit le même effet qu'une question sur ce nœud.
\end{enumerate}

\subsection{Exercice et preuve}

\begin{center}
\resizebox{0.97\textwidth}{!}{%
\begin{sequencediagram}
\newthread{u}{Utilisateur}
\newinst[1.2]{p}{Panneau}
\newinst[1.6]{e}{Moteur pédagogique}
\newinst[1.6]{v}{Vérification}
\newinst[1.6]{a}{Adaptateur}
\begin{callself}{e}{liaisons résolues, critère encore faux}{offre}
\end{callself}
\mess{e}{{exercise.update (consigne)}}{p}
\mess{u}{modifie la scène lui-même}{a}
\mess{a}{{\detokenize{event.state_changed (origin user)}}}{e}
\begin{call}{e}{vérifier les prédicats}{v}{preuve}
\begin{call}{v}{adapter.observe}{a}{observation}
\end{call}
\end{call}
\begin{callself}{e}{mise à jour de la maîtrise}{}
\end{callself}
\mess{e}{{evidence.added (verified), learner.update}}{p}
\mess{p}{vérifié dans l'application}{u}
\end{sequencediagram}}
\end{center}

Sans succès, l'utilisateur peut demander un indice (révélés dans l'ordre). Après \code{max\_attempts} échecs, Wysp propose la solution en mode guidé ou \code{do}. Si l'application ne permet pas de vérifier, la déclaration de l'utilisateur est enregistrée comme preuve \code{declared}, affichée comme non vérifiée.

\subsection{Développer avec Kiro}

\begin{figure}[h]
\centering
\begin{tikzpicture}[font=\sffamily\scriptsize, node distance=0.35cm and 0.45cm,
  st/.style={draw=#1, fill=#1!8, rounded corners=2pt, align=center, minimum height=0.95cm, text width=2.15cm},
  arr/.style={-{Stealth[length=2mm]}, thick, color=surface}]
  \node[st=pa] (spec) {Spec\\\tiny exigences, design, propriétés};
  \node[st=pa, right=of spec] (task) {Tâche\\\tiny « Start task »};
  \node[st=pumpkin, right=of task] (pre) {Hook\\\tiny task-preflight};
  \node[st=pb, right=of pre] (impl) {Implémentation\\\tiny agent + humain};
  \node[st=pumpkin, right=of impl] (save) {Hooks à la sauvegarde\\\tiny format, frontière, schémas, contrat};
  \node[st=pumpkin, below=0.9cm of save] (post) {Hook\\\tiny task-tests};
  \node[st=pc, left=of post] (guard) {Agent\\\tiny contract-guardian};
  \node[st=pb, left=of guard] (pr) {Pull request\\\tiny revue du responsable};
  \node[st=okgreen, left=of pr] (ci) {CI verte\\\tiny tests, kit, frontière};
  \node[st=okgreen, left=of ci] (merge) {Fusion\\\tiny au moins 1 fois par jour};
  \draw[arr] (spec) -- (task); \draw[arr] (task) -- (pre); \draw[arr] (pre) -- (impl); \draw[arr] (impl) -- (save);
  \draw[arr] (save) -- (post); \draw[arr] (post) -- (guard); \draw[arr] (guard) -- (pr); \draw[arr] (pr) -- (ci); \draw[arr] (ci) -- (merge);
  \draw[arr, dashed, color=muted] (save.north) |- ++(0,0.45) -| node[pos=0.25, above, text=muted] {échec : correction immédiate} (impl.north);
\end{tikzpicture}
\caption{Boucle de développement d'une tâche de spec.}
\end{figure}

\begin{enumerate}[leftmargin=1.2em]
  \item Chaque personne ouvre la spec de son périmètre dans Kiro et lance les tâches dans l'ordre (« Start task »). Le hook \code{task-preflight} rappelle les exigences, le responsable et le processus de changement de contrat.
  \item Pendant l'implémentation, les hooks de sauvegarde formatent le code et lancent les contrôles déterministes. Ils ne coûtent rien quand rien n'est concerné.
  \item À la fin de la tâche, \code{task-tests} exécute les tests des paquets modifiés. Les tâches de propriétés sont obligatoires.
  \item Avant la pull request, l'agent \code{contract-guardian} (lecture seule) produit un tableau de conformité. Le responsable du périmètre touché relit. La CI doit être verte.
  \item Le journal \file{docs/kiro-journal.md} note ce que Kiro a produit et ce qui a été corrigé à la main : c'est la matière de la section « How we used Kiro » de la soumission.
\end{enumerate}

\subsection{Changer le contrat}

\begin{enumerate}[leftmargin=1.2em]
  \item Écrire le besoin dans la spec concernée, puis modifier \file{.kiro/specs/adapter-contract/} (exigence et design).
  \item Modifier les modèles Pydantic dans \file{core/src/wysp_core/contract/} et augmenter \code{contract\_version} (mineure si ajout optionnel).
  \item Le hook \code{contract-schema-sync} régénère \file{schemas/} et les types TypeScript.
  \item Lancer \code{wysp-contract --offline} pour les deux tuteurs : les enregistrements montrent qui serait cassé.
  \item Obtenir l'accord de \PA{} et du responsable de chaque adaptateur touché. Après J0, pas de changement majeur sans accord des trois.
\end{enumerate}

\subsection{Créer un nouveau tuteur avec Wysp}

C'est le parcours que le jury doit voir comme valeur du squelette, et que le README présentera en quatre étapes.

\begin{enumerate}[leftmargin=1.2em]
  \item \textbf{Manifeste} : identifiant, niveau d'intégration, capacités, types d'entités, schéma de références.
  \item \textbf{Observation} : traduire l'état de l'application en entités (le SDK fournit client, empreinte et validation).
  \item \textbf{Actions} : un fichier par action avec schéma, risque, post-conditions et chemin d'interface ; le hook \code{new-action-integration} câble le reste.
  \item \textbf{Profil} : compétences, concepts, leçons et exercices en YAML, validés par \code{wysp validate-profile}.
\end{enumerate}
Puis \code{wysp-contract --tutor <nom>} indique si l'adaptateur est conforme. Aucun changement du cœur ni du panneau n'est nécessaire.

\section{Mise en place Kiro dans le dépôt}

Tout ce qui suit est déjà créé dans le dépôt, en anglais pour le jury. Le dossier \file{.kiro/} est à la racine et doit rester versionné.

\subsection{Steering : les règles que l'agent applique}

{\small\setlength{\tabcolsep}{4pt}
\begin{xltabular}{\textwidth}{@{}L{3.5cm}L{2.6cm}Y@{}}
\toprule
\textbf{Fichier} & \textbf{Inclusion} & \textbf{Contenu} \\
\midrule
\file{product.md} & toujours & Vision, vocabulaire (cœur, adaptateur, profil, tuteur, provenance\dots), promesse produit, règles non négociables, cible de la démo \\
\file{tech.md} & toujours & Pile technique, variables d'environnement, commandes (marquées « prévues » tant que la tâche qui les crée n'est pas faite), notes Windows \\
\file{structure.md} & toujours & Arborescence, responsables par dossier, règles de frontière, nommage, processus de changement de contrat, branches et messages de commit \\
\file{trust-and-safety.md} & toujours & Provenance, niveaux de risque, permissions, annulation, sécurité locale, vie privée, audit \\
\file{testing.md} & toujours & Couches de tests, identifiants des propriétés (\code{CON-P}, \code{CORE-P}\dots), règles (pas de vrai LLM en test, pas de contrat simulé) \\
\file{adapter-contract.md} & fichiers \file{apps/}, \file{sdk/}, contrat, passerelle & Guide court du contrat pour écrire un adaptateur, règles propres à Blender et à Chrome \\
\file{pedagogy.md} & profils, pédagogie, vérification, orchestrateur & Stratégie de tutorat, format des profils, exercices, preuves, maîtrise \\
\file{panel-ui.md} & \file{panel/}, tableau de bord & Modes d'affichage, identité visuelle, langage de provenance, règles d'interaction, accessibilité \\
\file{demo-check.md} & manuelle (\code{/demo-check}) & Liste de contrôle avant tournage ; remplace les hooks manuels absents de Kiro IDE 1.x \\
\file{release-checklist.md} & manuelle & Liste de soumission Kiroween 2026 \\
\bottomrule
\end{xltabular}}

\subsection{Specs : six périmètres, une traçabilité complète}

Chaque spec suit le format du lauréat 2025 : exigences en EARS, design avec propriétés de correction reliées aux exigences, tâches reliées aux exigences. Nous ajoutons pour chaque tâche un \textbf{responsable} et un \textbf{jalon}, et les tâches de propriétés sont obligatoires.

{\small\setlength{\tabcolsep}{4pt}
\begin{xltabular}{\textwidth}{@{}L{4.3cm}L{1.3cm}r r r r@{}}
\toprule
\textbf{Spec} & \textbf{Pilote} & \textbf{Exig.} & \textbf{Critères} & \textbf{Propr.} & \textbf{Tâches (sous-t.)} \\
\midrule
\file{adapter-contract} & \PA & 12 & 53 & 11 & 13 (7) \\
\file{core-orchestrator} & \PA & 11 & 47 & 12 & 14 (6) \\
\file{pedagogy-verification} & \PA & 6 & 26 & 9 & 9 (4) \\
\file{tutor-panel} & \PB & 6 & 28 & 6 & 11 (4) \\
\file{blender-tutor} & \PC & 8 & 31 & 5 & 12 (4) \\
\file{web-tutor} & \PC{} (\PB) & 9 & 32 & 6 & 12 (4) \\
\midrule
\textbf{Total} & & \textbf{52} & \textbf{217} & \textbf{49} & \textbf{71 (29)} \\
\bottomrule
\end{xltabular}}

Trois tâches seulement sont optionnelles (marquées \code{*}) : un adaptateur de niveau visuel, un fournisseur Bedrock et le pont MCP de développement. Elles ne sont commencées que si le planning le permet.

\subsection{Hooks : contrôles automatiques}

Les hooks sont au format actuel \file{.kiro/hooks/*.json} (version v1), lu par Kiro IDE 1.x et Kiro CLI. Les contrôles fréquents sont des \textbf{commandes} déterministes qui appellent \file{scripts/kiro/hook.py} (Python standard, testé) : elles ne consomment pas de crédits et ne font rien tant que la tâche qui crée leur cible n'est pas terminée. Les prompts d'agent sont réservés à la revue de sécurité et à l'intégration d'une nouvelle action.

{\small\setlength{\tabcolsep}{4pt}
\begin{xltabular}{\textwidth}{@{}L{4.2cm}L{2.4cm}L{1.4cm}Y@{}}
\toprule
\textbf{Hook} & \textbf{Déclencheur} & \textbf{Type} & \textbf{Effet} \\
\midrule
\code{core-boundary-guard} & sauvegarde \file{core/*.py} & commande & Refuse les imports d'hôte ou d'adaptateur et le vocabulaire d'application dans le cœur \\
\code{trust-review} & sauvegarde \file{trust/}, \file{actions/}, manifestes & agent & Revue contre les règles de confiance : risque, annulation, pas d'évaluation de code, pas de fuite de secret, permissions minimales \\
\code{task-preflight} & avant une tâche & agent & Rappelle exigences, responsable, dépendances et processus de contrat \\
\code{contract-schema-sync} & sauvegarde du contrat & commande & Régénère JSON Schema et types TypeScript \\
\code{adapter-contract-offline} & sauvegarde d'un adaptateur & commande & Kit de contrat hors ligne pour le tuteur touché \\
\code{new-action-integration} & création d'un fichier d'action & agent & Câble la nouvelle action : registre, annulation, fixture du kit, README (équivalent de leurs hooks de gestes) \\
\code{format-python}, \code{format-typescript} & sauvegarde & commande & ruff et prettier sur les fichiers concernés \\
\code{profile-validate} & sauvegarde d'un profil & commande & \code{wysp validate-profile} \\
\code{task-tests} & après une tâche & commande & Tests des paquets modifiés (pytest, Vitest) \\
\bottomrule
\end{xltabular}}

\begin{notebox}
\small Les hooks de fichiers ne se déclenchent que pour les modifications faites par l'agent, pas pour une sauvegarde manuelle dans l'éditeur. La CI exécute donc les mêmes contrôles (\code{python scripts/kiro/hook.py boundaries}, tests, export des schémas) pour couvrir le travail fait à la main.
\end{notebox}

\subsection{Agents spécialisés et MCP}

{\small\setlength{\tabcolsep}{4pt}
\begin{xltabular}{\textwidth}{@{}L{3.4cm}Y L{3.6cm}@{}}
\toprule
\textbf{Agent} & \textbf{Rôle} & \textbf{Outils} \\
\midrule
\code{blender-adapter} & Implémente et teste l'adaptateur Blender et le SDK Python, avec les règles de fil principal, d'annulation et d'événements & lecture, écriture, shell, web, Context7 \\
\code{browser-adapter} & Implémente et teste l'extension Chrome et le SDK TypeScript, avec les règles MV3, de localisateurs et de vie privée & lecture, écriture, shell, web, Context7, Playwright \\
\code{contract-guardian} & Relecteur en lecture seule : frontière, contrat, actions, confiance, tests ; rend un tableau de conformité & lecture, shell limité (\code{git diff}, \code{uv run}) \\
\bottomrule
\end{xltabular}}

Configuration MCP proposée dans \file{docs/kiro/mcp.example.json}, à copier dans \file{.kiro/settings/mcp.json} (l'agent n'a pas le droit d'écrire dans ce dossier) :
\begin{itemize}[leftmargin=1.2em]
  \item \textbf{Context7} : documentation à jour de \code{bpy}, React Flow, Chrome Extensions, FastAPI.
  \item \textbf{Playwright MCP} : piloter Chrome avec l'extension et le tableau de bord pendant le développement.
  \item \textbf{wysp-dev} (désactivé jusqu'à la tâche optionnelle 13 de \file{core-orchestrator}) : pont MCP en lecture seule vers un tuteur en cours d'exécution (adaptateurs, observation, tableau, audit). Kiro peut alors inspecter l'état réel de Blender pendant qu'on développe Blender.
\end{itemize}

\subsection{Utilisation au quotidien}

\begin{itemize}[leftmargin=1.2em]
  \item Ouvrir la spec de son périmètre, lancer la prochaine tâche dont on est responsable, laisser les hooks tourner, relire le diff.
  \item Référencer une spec dans le chat avec \code{\#spec:adapter-contract} ; inclure une règle manuelle avec \code{\#demo-check}.
  \item Déléguer à un agent : « utilise l'agent \code{blender-adapter} pour la tâche 4 de \file{blender-tutor} ».
  \item Avant chaque pull request : « lance \code{contract-guardian} sur ma branche ».
  \item Noter dans \file{docs/kiro-journal.md} les moments où Kiro a été décisif (spec générée, bug trouvé par un hook, propriété qui a cassé) avec la date : c'est la matière du texte de soumission et de la vidéo.
\end{itemize}

\section{Plan d'implémentation}

\subsection{Hypothèses de capacité}

\begin{itemize}[leftmargin=1.2em]
  \item 25 jours calendaires du lundi 9 novembre au vendredi 4 décembre. Nous comptons 22 jours de développement par personne et gardons les 2 et 3 décembre pour la vidéo et la soumission.
  \item Les charges ci-dessous sont en jours-personne. L'estimation « avec Kiro » applique un facteur 0,65 à l'estimation brute, parce que l'agent écrit une grande partie du code à partir des specs et que l'humain relit, teste et corrige. \textbf{Ce facteur est une hypothèse} : à recalibrer au jalon J0 en comparant le temps réel des premières tâches à leur estimation.
  \item Capacité : 3 $\times$ 22 = 66 jours-personne. Besoin estimé : 56 avec Kiro (86 brut). Marge d'environ 15\,\%, qui sert aux imprévus d'intégration.
\end{itemize}

\subsection{Jalons et critères de sortie}

\begin{xltabular}{\textwidth}{@{}L{2.6cm}Y@{}}
\toprule
\textbf{Jalon} & \textbf{Critère de sortie (vérifié ensemble, avec \code{/demo-check})} \\
\midrule
\textbf{J0} mer.\ 11 nov. & Contrat 1.0.0 figé (modèles, schémas, types TS), adaptateur factice connecté à \code{wysp serve --mock}, panneau « hello » qui liste les adaptateurs, CI verte avec contrôle de frontière. \\
\textbf{J1} jeu.\ 19 nov. & Blender de bout en bout \textbf{en mode texte} : expliquer la pile de modificateurs de la citrouille, tableau en disposition \code{stack}, proposition \code{add\_modifier}, approbation, exécution, vérification, annulation. Tableau de bord de démo v1. \\
\textbf{J2} ven.\ 27 nov. & Navigateur de bout en bout sur \textbf{le même cœur} (diff de \file{core/} vide pour ce travail) : expliquer le filtre de région, création guidée d'un rapport vérifiée. Voix avec interruption. Exercices et preuves sur les deux applications. \\
\textbf{J3} mar.\ 1\textsuperscript{er} déc. & Gel des fonctionnalités. Kit de contrat vert pour les trois adaptateurs, README et guide d'adaptateur, scénario de démo répété deux fois sans incident. \\
\textbf{Soumission} ven.\ 4 déc. & Formulaire Devpost envoyé avant 12\,h GMT (limite officielle 22\,h GMT). \\
\bottomrule
\end{xltabular}

\subsection{Calendrier}

\begin{center}
\resizebox{\textwidth}{!}{%
\begin{ganttchart}[
    time slot format=isodate,
    x unit=0.56cm, y unit chart=0.52cm, y unit title=0.55cm,
    hgrid={*1{draw=muted!25}}, vgrid={*1{draw=muted!15}},
    title/.style={draw=none, fill=night!8},
    title label font=\sffamily\tiny,
    bar label font=\sffamily\scriptsize, group label font=\sffamily\scriptsize\bfseries,
    milestone label font=\sffamily\scriptsize\bfseries,
    bar height=0.55, group top shift=0.25, group height=0.25, group peaks height=0.12,
    bar/.append style={draw=none, rounded corners=1pt},
    milestone/.append style={fill=danger, draw=none},
    link/.style={-{Stealth[length=1.5mm]}, color=muted}
  ]{2026-11-09}{2026-12-04}
  \gantttitle{novembre}{22} \gantttitle{déc.}{4} \\
  \gantttitlecalendar{day} \\
  \ganttgroup[group/.append style={fill=pa}]{A, cœur et contrat}{2026-11-09}{2026-12-01} \\
  \ganttbar[bar/.append style={fill=pa!45}]{Contrat, factice, squelette}{2026-11-09}{2026-11-11} \\
  \ganttbar[bar/.append style={fill=pa!45}]{Passerelle, LLM, tableau}{2026-11-12}{2026-11-15} \\
  \ganttbar[bar/.append style={fill=pa!45}]{Porte, orchestrateur, kit}{2026-11-16}{2026-11-19} \\
  \ganttbar[bar/.append style={fill=pa!45}]{Exercices, leçons, apprenant}{2026-11-20}{2026-11-27} \\
  \ganttbar[bar/.append style={fill=pa!25}]{Intégration, robustesse}{2026-11-28}{2026-12-01} \\
  \ganttgroup[group/.append style={fill=pb}]{B, expérience}{2026-11-09}{2026-12-03} \\
  \ganttbar[bar/.append style={fill=pb!40}]{Panneau, ui\_api}{2026-11-09}{2026-11-13} \\
  \ganttbar[bar/.append style={fill=pb!40}]{Tableau de bord de démo}{2026-11-13}{2026-11-16} \\
  \ganttbar[bar/.append style={fill=pb!40}]{Rendu du tableau, sous-titres}{2026-11-14}{2026-11-19} \\
  \ganttbar[bar/.append style={fill=pb!40}]{Contenu des profils}{2026-11-20}{2026-11-25} \\
  \ganttbar[bar/.append style={fill=pb!40}]{Voix, synthèse, exercices}{2026-11-21}{2026-11-27} \\
  \ganttbar[bar/.append style={fill=pb!25}]{Accessibilité, vidéo}{2026-11-28}{2026-12-03} \\
  \ganttgroup[group/.append style={fill=pc}]{C, adaptateurs}{2026-11-09}{2026-12-01} \\
  \ganttbar[bar/.append style={fill=pc!40}]{SDK Python, extension Blender}{2026-11-09}{2026-11-13} \\
  \ganttbar[bar/.append style={fill=pc!40}]{Observation, actions, événements}{2026-11-14}{2026-11-19} \\
  \ganttbar[bar/.append style={fill=pc!40}]{SDK TS, surbrillance, kit Blender}{2026-11-20}{2026-11-22} \\
  \ganttbar[bar/.append style={fill=pc!40}]{Extension Chrome}{2026-11-23}{2026-11-27} \\
  \ganttbar[bar/.append style={fill=pc!25}]{README, guide, tests live}{2026-11-28}{2026-12-01} \\
  \ganttbar[bar/.append style={fill=okgreen!45}]{Tournage, montage, envoi}{2026-12-02}{2026-12-04} \\
  \ganttmilestone{Jalons J0, J1, J2, J3, envoi}{2026-11-11} \ganttmilestone{}{2026-11-19} \ganttmilestone{}{2026-11-27} \ganttmilestone{}{2026-12-01} \ganttmilestone{}{2026-12-04}
\end{ganttchart}}
\end{center}

\subsection{Les trois sections de travail}

\subsubsection*{\textcolor{pa}{Personne A : cœur et contrat}}

\begin{xltabular}{\textwidth}{@{}L{3.2cm}Y@{}}
\toprule
Mission & Définir et tenir le contrat ; construire le cœur qui rend les deux tuteurs possibles sans code d'application. \\
Dossiers & \file{core/}, \file{schemas/}, \file{contract-tests/} (kit et adaptateur factice). \\
Specs et tâches & \file{adapter-contract} 1--6, 9, 10, 13 ; \file{core-orchestrator} 1--3, 5--8, 10, 12, 14 (et 5.1, 13 si le temps le permet) ; \file{pedagogy-verification} 1--9. \\
Fournit à & \PB{} : protocole du panneau (types générés), fixtures de correctifs du tableau, LLM scénarisé. \PC{} : contrat figé à J0, kit en ligne lundi 16 novembre, chargeur de profils lundi 16 novembre. \\
Dépend de & \PB{} pour \code{ui\_api} (tâche 4) ; \PC{} pour les retours sur le contrat avant J0. \\
Charge & 32 jours bruts, environ 21 avec Kiro. C'est la charge la plus forte, d'où les leviers ci-dessous. \\
Ligne de coupe & Dans l'ordre : pont MCP et Bedrock (optionnels), recommandations (remplacées par un ordre fixe des leçons), mode \code{--offline} du kit, injection de pannes en CI, résumé glissant de la conversation. \\
\bottomrule
\end{xltabular}

\subsubsection*{\textcolor{pb}{Personne B : expérience}}

\begin{xltabular}{\textwidth}{@{}L{3.2cm}Y@{}}
\toprule
Mission & Rendre le tutorat visible, audible et agréable : panneau, tableau, voix, tableau de bord de démo, contenu pédagogique, vidéo. \\
Dossiers & \file{panel/}, \file{core/src/wysp_core/ui_api/}, \file{core/src/wysp_core/narration/}, \file{apps/web-tutor/demo-dashboard/}, \file{apps/*/profile/}. \\
Specs et tâches & \file{tutor-panel} 1--11 ; \file{core-orchestrator} 4, 9, 11 ; \file{web-tutor} 1, 9 ; \file{blender-tutor} 9 ; \file{adapter-contract} 1 (partie npm). \\
Fournit à & \PA{} : serveur \code{ui\_api}, narration. \PC{} : tableau de bord accessible lundi 16 novembre (cible de l'instantané), style de l'overlay. Tous : identité visuelle, vidéo, visuels du README. \\
Dépend de & \PA{} : types du protocole (J0), fixtures du tableau (tâche 6), chargeur de profils (lundi 16 novembre). \PC{} : relecture technique des profils. \\
Charge & 27 jours bruts, environ 17,5 avec Kiro. \\
Ligne de coupe & Synthèse vocale serveur (garder celle du navigateur), dispositions \code{graph} et \code{sequence}, traduction française complète, interruption par détection d'énergie (garder Échap et l'appui pour parler). \\
\bottomrule
\end{xltabular}

\subsubsection*{\textcolor{pc}{Personne C : adaptateurs}}

\begin{xltabular}{\textwidth}{@{}L{3.2cm}Y@{}}
\toprule
Mission & Brancher deux applications très différentes sur le même contrat, sans jamais toucher au cœur. \\
Dossiers & \file{sdk/}, \file{apps/blender-tutor/adapter/} et \file{demo/}, \file{apps/web-tutor/extension/}. \\
Specs et tâches & \file{adapter-contract} 7, 8, 11 ; \file{blender-tutor} 1--8, 10--12 ; \file{web-tutor} 2--8, 10--12. \\
Fournit à & \PA{} : retours sur le contrat avant J0, fixtures de contrat et enregistrements. \PB{} : scène de démo, catalogue d'actions pour les profils. \\
Dépend de & \PA{} : contrat (J0), kit (lundi 16 novembre). \PB{} : tableau de bord de démo (lundi 16 novembre). \\
Charge & 27,5 jours bruts, environ 18 avec Kiro. \\
Ligne de coupe & \code{apply\_modifier}, \code{assign\_material}, \code{click\_element}, \code{follow\_link} ; mise en évidence des modificateurs via l'éditeur de propriétés ; traversée des racines Shadow DOM. \\
\bottomrule
\end{xltabular}

\begin{keybox}[Leviers de rééquilibrage, décidés au point du jalon]
Si A a plus d'un jour de retard à J1 : C reprend le lanceur Playwright et le mode hors ligne du kit, B reprend le moteur de recommandation et les tests de scénarios des leçons. Si C est en retard à J2 : A aide sur les actions web (après J2, la charge de A baisse), et l'adaptateur web se limite à \code{select\_option}, \code{fill\_field}, \code{set\_checked}, \code{focus\_element}.
\end{keybox}

\subsection{Planning semaine par semaine}

{\small
\begin{xltabular}{\textwidth}{@{}L{1.9cm}Y Y Y@{}}
\toprule
\textbf{Semaine} & \textcolor{pa}{\textbf{A}} & \textcolor{pb}{\textbf{B}} & \textcolor{pc}{\textbf{C}} \\
\midrule
9--15 nov. & Monorepo, modèles et export, prédicats et empreinte, adaptateur factice, squelette du cœur, passerelle, session (\textbf{J0}) ; puis LLM scénarisé et fournisseur réel, modèle du tableau. & Squelette du panneau et client du protocole (\textbf{J0}), serveur \code{ui\_api}, identité visuelle et modes, magasin du tableau ; début du tableau de bord. & Scène de démo, SDK Python (jeudi 12), extension Blender et connexion (jeudi 12, vendredi 13), début de l'observation. \\
16--22 nov. & Kit en ligne et chargeur de profils (lundi 16), porte et chaîne d'action, orchestrateur et outils, service de vérification (\textbf{J1}) ; puis tests de scénarios, début des exercices. & Tableau de bord v1 (lundi 16), rendu du tableau (\code{stack}, \code{tree}), sous-titres avec synthèse du navigateur, dialogue de permission (\textbf{J1}) ; puis narration côté cœur et profil Blender. & Observation (lundi 16), actions et annulation, événements et origine, fixture du kit (\textbf{J1}) ; puis SDK TypeScript, surbrillance, kit Blender en mode sans interface. \\
23--29 nov. & Moteur d'exercices, déroulé des leçons, modèle de l'apprenant, comportements sur événements ; appui à l'intégration web (\textbf{J2}). & Profil web, voix (transcription, interruption), synthèse serveur, carte d'exercice et progression, dispositions \code{graph} et \code{sequence} (\textbf{J2}). & Extension Chrome : squelette, instantané et masquage, localisateurs, actions, événements, overlay, fixture et lanceur (\textbf{J2}). \\
30 nov.--4 déc. & Intégration et corrections (lundi 30), gel (\textbf{J3}), robustesse, aide au tournage avec le LLM scénarisé en secours. & Accessibilité, finitions, script et tournage (1\textsuperscript{er}--2 déc.), montage (3 déc.), visuels du README. & README des tuteurs, guide d'adaptateur, tests live ; opère Blender et Chrome pendant le tournage. \\
\bottomrule
\end{xltabular}}

\subsection{Synchronisation et règles d'équipe}

\begin{itemize}[leftmargin=1.2em]
  \item \textbf{Point quotidien de 15 minutes} à heure fixe : fait hier, prévu aujourd'hui, blocage. Les blocages de contrat passent avant tout le reste.
  \item \textbf{Séance d'intégration à chaque jalon} (2 heures, les trois ensemble) : on exécute \code{/demo-check}, on note les écarts, on décide des coupes.
  \item \textbf{Contrat figé à J0.} Ensuite, seuls les ajouts optionnels (version mineure) passent, selon le processus de la section 4.6.
  \item \textbf{Branches courtes} \code{a/…}, \code{b/…}, \code{c/…}, fusionnées au moins une fois par jour. Deux personnes travaillent sur \file{core-orchestrator/tasks.md} (A et B) : fusionner souvent pour éviter les conflits sur les cases cochées.
  \item \textbf{Une session Kiro par personne et par spec.} Ne pas lancer « Run all tasks » sur une spec partagée : chacun lance ses propres tâches.
\end{itemize}

\subsection{Définition de « fini » pour une tâche}

\begin{enumerate}[leftmargin=1.2em]
  \item Le code fait ce que disent les critères d'acceptation cités par la tâche.
  \item Les tests de la tâche passent, y compris les tests de propriétés obligatoires ; aucun test n'appelle un vrai LLM ou un vrai service vocal.
  \item Contrôle de frontière propre ; kit de contrat vert pour le tuteur touché (si concerné).
  \item Documentation mise à jour (tableau des actions du README, commandes de \file{tech.md} si une commande « prévue » devient réelle).
  \item Relecture par le responsable du périmètre, CI verte, fusion sur \code{main}, case cochée dans \file{tasks.md}.
\end{enumerate}

\subsection{Priorités (MoSCoW)}

\begin{xltabular}{\textwidth}{@{}L{2.4cm}Y@{}}
\toprule
\textbf{Priorité} & \textbf{Contenu} \\
\midrule
Indispensable & Contrat et kit ; cœur avec outils, tableau, provenance, porte, vérification et annulation ; mode guidé ; narration avec interruption ; un exercice vérifié par application ; Blender avec \code{select\_object}, \code{show\_properties\_tab}, \code{add\_modifier}, \code{set\_modifier\_property}, \code{remove\_modifier} ; web avec \code{focus\_element}, \code{select\_option}, \code{fill\_field}, \code{set\_checked} ; tableau de bord ; panneau en modes \code{side} et \code{compact} ; appui pour parler ; synthèse du navigateur. \\
Important & Déroulé complet des leçons, modèle de l'apprenant, dispositions \code{graph} et \code{sequence}, synthèse vocale serveur, interruption par la voix, actions restantes, mode hors ligne du kit, français complet. \\
Souhaitable & Recommandations, injection de pannes en CI, pont MCP, Bedrock, mode \code{full} du panneau. \\
Exclu du prototype & Adaptateur de niveau visuel (documenté comme extension possible), multi-utilisateur, déploiement distant. \\
\bottomrule
\end{xltabular}

\section{Démonstration, risques et questions ouvertes}

\subsection{Scénario de la vidéo (moins de 3 minutes)}

La vidéo doit montrer un produit utilisable, la réutilisation réelle du cœur et l'usage de Kiro. Elle est tournée en anglais (ou sous-titrée en anglais), avec le LLM réel ; le LLM scénarisé reste prêt en secours.

\begin{xltabular}{\textwidth}{@{}L{1.6cm}L{3.0cm}Y@{}}
\toprule
\textbf{Temps} & \textbf{Écran} & \textbf{Ce qu'on montre} \\
\midrule
0:00--0:15 & Titre, Blender & Le problème : apprendre un logiciel sur son propre travail. Wysp en une phrase. \\
0:15--0:50 & Blender + panneau & « Explain this pumpkin's modifiers. » Le tableau se construit en pile (nœuds observés), la narration met en évidence Subdivision puis Solidify dans Blender. \\
0:50--1:05 & idem & Interruption à la voix : « Wait, why is Solidify after Subdivision? » Réponse avec un nœud \emph{illustratif} « if Solidify came first », marqué « for example ». \\
1:05--1:35 & idem & « Round the candle's edges. » Dialogue de permission, approbation, ajout du Bevel, badge « verified in Blender ». Puis l'exercice « 3 segments » fait par l'utilisateur lui-même et vérifié. \\
1:35--2:15 & Chrome + panneau & Même panneau, même voix. « What does the region filter do? » Le tableau montre les filtres et le tableau des commandes. Création guidée du rapport : surbrillance des contrôles, l'utilisateur agit, rapport vérifié dans la liste. \\
2:15--2:40 & Code, terminal & Seul l'adaptateur change : arborescence \file{apps/}, \code{git diff} vide sur \file{core/}, le même kit de contrat vert sur factice, Blender et Chrome, le hook de frontière en action. \\
2:40--2:55 & \file{.kiro/} & Specs avec propriétés, hooks, agents ; « build your own tutor in 4 steps ». \\
2:55--3:00 & Logo & Lien du dépôt. \\
\bottomrule
\end{xltabular}

\subsection{Risques et mitigations}

\begin{xltabular}{\textwidth}{@{}L{4.3cm}L{1.3cm}Y@{}}
\toprule
\textbf{Risque} & \textbf{Niveau} & \textbf{Mitigation} \\
\midrule
Glissement du contrat après J0, intégration tardive & élevé & Contrat figé à J0, kit en ligne lundi 16 novembre, enregistrements hors ligne, intégration de bout en bout dès J1 avec l'adaptateur factice puis Blender. \\
Charge de A la plus forte & élevé & B porte \code{ui\_api} et la narration ; lignes de coupe définies ; leviers de rééquilibrage à J1 et J2. \\
Fiabilité du LLM (outils mal appelés, latence) pendant la démo & élevé & Outils peu nombreux et typés, erreurs structurées, 8 tours maximum, scénario rejoué par \code{ScriptedLLM} en secours, prompts testés sur les deux scénarios. \\
Voix : écho, faux déclenchements, permission micro & moyen & Appui pour parler par défaut, Échap toujours actif, annulation d'écho, repli sur le texte et sur la synthèse du navigateur. \\
Threads Blender (plantage si \code{bpy} hors du fil principal) & moyen & File et pompe par minuteur, propriété BLD-P1, tests avec \code{bpy} de PyPI. \\
Service worker MV3 arrêté par Chrome & moyen & Battement de 20\,s, reconnexion et état dans \code{chrome.storage.session}. \\
Changements d'API entre versions de Blender & moyen & Cible unique Blender 5.2 LTS, Context7, extraction RNA générique plutôt que code par type. \\
Coût des appels LLM et vocaux & faible & Fournisseur compatible OpenAI configurable (y compris modèle local), cache de la synthèse vocale. \\
Règles du concours mal lues (page non accessible automatiquement) & moyen & Relecture humaine des règles officielles cette semaine ; liste \code{/release-checklist}. \\
Dépassement de la durée de la vidéo & faible & Script minuté ci-dessus, répétition à J3. \\
\bottomrule
\end{xltabular}

\subsection{Questions ouvertes à trancher cette semaine}

\begin{enumerate}[leftmargin=1.2em]
  \item \textbf{Fournisseur LLM} : quel modèle avec appels d'outils, quel budget ? Une offre compatible OpenAI (ou Bedrock si nous avons des crédits AWS) est prévue par l'architecture.
  \item \textbf{Synthèse vocale} : service externe (meilleure voix) ou synthèse du navigateur seule (aucun coût) ?
  \item \textbf{Répartition A, B, C} : la proposition suppose une personne à l'aise en Python asynchrone (A), une en React et design (B), une curieuse de Blender et des extensions Chrome (C). À ajuster selon les profils réels.
  \item \textbf{Machines} : Blender 5.2 LTS et Chrome récent sur les trois postes ; une machine de tournage avec micro correct.
  \item \textbf{Langue} : dépôt et vidéo en anglais (recommandé pour le jury), panneau bilingue.
  \item \textbf{Licences} : MIT à la racine, GPL-3.0-or-later pour l'extension Blender. À confirmer par l'équipe.
\end{enumerate}

\appendix
\section{Annexes}

\subsection{Fichiers créés dans le dépôt}

\begin{xltabular}{\textwidth}{@{}L{6.6cm}Y@{}}
\toprule
\textbf{Chemin} & \textbf{Contenu} \\
\midrule
\file{.kiro/steering/*.md} (10) & Règles de l'agent : produit, technique, structure, confiance, tests, contrat, pédagogie, panneau, deux listes manuelles \\
\file{.kiro/specs/*/} (6 specs) & \file{requirements.md}, \file{design.md}, \file{tasks.md} pour \file{adapter-contract}, \file{core-orchestrator}, \file{pedagogy-verification}, \file{tutor-panel}, \file{blender-tutor}, \file{web-tutor} \\
\file{.kiro/hooks/wysp-guards.json} & \code{core-boundary-guard}, \code{trust-review}, \code{task-preflight} \\
\file{.kiro/hooks/wysp-contract.json} & \code{contract-schema-sync}, \code{adapter-contract-offline}, \code{new-action-integration} \\
\file{.kiro/hooks/wysp-quality.json} & \code{format-python}, \code{format-typescript}, \code{profile-validate}, \code{task-tests} \\
\file{.kiro/agents/*.md} (3) & \code{blender-adapter}, \code{browser-adapter}, \code{contract-guardian} \\
\file{docs/kiro/mcp.example.json} & Configuration MCP proposée (Context7, Playwright, pont \code{wysp-dev} désactivé) \\
\file{scripts/kiro/hook.py} & Contrôles appelés par les hooks et la CI (frontière, schémas, contrat, profils, format, tests) \\
\file{docs/architecture/} & Ce document (source LaTeX, sections séparées) \\
\file{docs/notes/working-notes.md} & Notes de travail : faits vérifiés, décisions, jalons \\
\file{README.md}, \file{LICENSE}, \file{.gitignore} & Présentation du dépôt, licence MIT, exclusions (en gardant \file{.kiro/}) \\
\bottomrule
\end{xltabular}

\subsection{Glossaire}

\begin{xltabular}{\textwidth}{@{}L{3.6cm}Y@{}}
\toprule
\textbf{Terme} & \textbf{Définition} \\
\midrule
Wysp Core (cœur) & Moteur Python commun à tous les tuteurs ; ne connaît aucune application. \\
Adaptateur & Composant propre à une application qui implémente le contrat. \\
Profil & Données pédagogiques (YAML) d'une application : compétences, concepts, leçons, exercices. \\
Tuteur & Adaptateur + profil, branchés sur le cœur ; un dossier par tuteur dans \file{apps/}. \\
Panneau & Interface React unique, commune aux tuteurs. \\
Tableau & Vue interactive du modèle visuel (nœuds, arêtes, étapes, annotations). \\
Provenance & Statut d'un fait : \pobs{} (observé), \pinf{} (déduit), \pill{} (exemple pédagogique). \\
Niveau d'intégration & \code{native} (API), \code{interface} (DOM, accessibilité), \code{visual} (capture, vision). \\
Porte de permission & Composant du cœur qui décide si une action proposée peut s'exécuter. \\
Preuve & Trace d'une vérification : prédicat, valeur réelle, valeur attendue, observation, date. \\
Empreinte d'état & Hachage canonique d'une observation, utilisé pour l'annulation et les appels invalides. \\
Kit de contrat & Suite de tests qui joue le rôle du cœur contre n'importe quel adaptateur. \\
EARS & Easy Approach to Requirements Syntax : « WHEN \dots{} THE SYSTEM SHALL \dots{} ». \\
Propriété de correction & Énoncé « pour toute entrée valide \dots{} », testé par génération aléatoire (Hypothesis, fast-check). \\
\bottomrule
\end{xltabular}

\subsection{Sources}

\begin{itemize}[leftmargin=1.2em]
  \item Note de concept « Project K » (fichier \file{project_k_rapport.pdf} du dépôt), renommée Wysp.
  \item Kiro InteractionKit, dépôt public : \url{https://github.com/Harpita-P/Kiro-InteractionKit} (cloné le 9 novembre 2026, dernier commit du 5 décembre 2025).
  \item Kiroween 2026, page du concours : \url{https://kiroween.devpost.com/} (période de soumission du 30 octobre au 4 décembre 2026 ; la page des règles a renvoyé une erreur 403 à la lecture automatique, à relire manuellement).
  \item Kiro Weekly Reading List, 1\textsuperscript{er} novembre 2026 (AWS Builder Center) : clôture le 4 décembre à 14\,h heure du Pacifique.
  \item Documentation Kiro (steering, specs, hooks, agents, MCP, skills) : \url{https://kiro.dev/docs/}.
  \item Blender 5.2 LTS, sortie le 14 juillet 2026, Python 3.13 ; \code{bpy} sur PyPI : \url{https://pypi.org/project/bpy/}.
  \item Chrome, WebSockets dans les service workers (Chrome 116 et suivants) : \url{https://developer.chrome.com/docs/extensions/how-to/web-platform/websockets}.
  \item Extensions Blender et licence GPL-3.0-or-later : \url{https://docs.blender.org/manual/en/latest/advanced/extensions/getting_started.html}.
\end{itemize}


\end{document}
